\subsection*{Textbook Formal Inventory}
This subsection records the exact formal material in \file{HAETAE_Security.ec}.  It is generated from the proof-surface EasyCrypt file, so the line numbers and statements are meant to be read next to the checked source.

\noindent\textbf{Chapter role.} top-level correctness and EUF-CMA theorem composition.

\noindent\textbf{Inventory size.} 65 context declarations, 1 definitions/game objects, 40 lemmas or relational theorems, and 1 explicit Hoare/Phoare judgments were found.

\subsubsection*{Theory Context, Imports, and Quantified Modules}
\paragraph{Context 1 (line 1)}
\noindent\textbf{EasyCrypt name.}
\begin{lstlisting}[language=EasyCrypt,basicstyle=\ttfamily\scriptsize]
imported theories
\end{lstlisting}
\noindent\textbf{Checked EasyCrypt statement.}
\begin{lstlisting}[language=EasyCrypt,basicstyle=\ttfamily\scriptsize]
require import AllCore Real Distr StdOrder.
\end{lstlisting}
\noindent\textbf{Formal mathematical reading.}
Mathematical context. This line imports or specializes earlier theories, making their carrier sets, functions, modules, and theorems available to the current file.

\noindent\textbf{Role in the proof.}
Use in this chapter. It fixes the proof context, imported mathematical language, or quantified module parameters for the following statements.

\paragraph{Context 2 (line 2)}
\noindent\textbf{EasyCrypt name.}
\begin{lstlisting}[language=EasyCrypt,basicstyle=\ttfamily\scriptsize]
imported theories
\end{lstlisting}
\noindent\textbf{Checked EasyCrypt statement.}
\begin{lstlisting}[language=EasyCrypt,basicstyle=\ttfamily\scriptsize]
require import Sig_ROM HAETAE_Scheme HAETAE_Algebra HAETAE_Distributions.
\end{lstlisting}
\noindent\textbf{Formal mathematical reading.}
Mathematical context. This line imports or specializes earlier theories, making their carrier sets, functions, modules, and theorems available to the current file.

\noindent\textbf{Role in the proof.}
Use in this chapter. It fixes the proof context, imported mathematical language, or quantified module parameters for the following statements.

\paragraph{Context 3 (line 3)}
\noindent\textbf{EasyCrypt name.}
\begin{lstlisting}[language=EasyCrypt,basicstyle=\ttfamily\scriptsize]
imported theories
\end{lstlisting}
\noindent\textbf{Checked EasyCrypt statement.}
\begin{lstlisting}[language=EasyCrypt,basicstyle=\ttfamily\scriptsize]
require import HAETAE_Reductions HAETAE_Assumptions.
\end{lstlisting}
\noindent\textbf{Formal mathematical reading.}
Mathematical context. This line imports or specializes earlier theories, making their carrier sets, functions, modules, and theorems available to the current file.

\noindent\textbf{Role in the proof.}
Use in this chapter. It fixes the proof context, imported mathematical language, or quantified module parameters for the following statements.

\paragraph{Context 4 (line 4)}
\noindent\textbf{EasyCrypt name.}
\begin{lstlisting}[language=EasyCrypt,basicstyle=\ttfamily\scriptsize]
imported theories
\end{lstlisting}
\noindent\textbf{Checked EasyCrypt statement.}
\begin{lstlisting}[language=EasyCrypt,basicstyle=\ttfamily\scriptsize]
require import HAETAE_Rejection.
\end{lstlisting}
\noindent\textbf{Formal mathematical reading.}
Mathematical context. This line imports or specializes earlier theories, making their carrier sets, functions, modules, and theorems available to the current file.

\noindent\textbf{Role in the proof.}
Use in this chapter. It fixes the proof context, imported mathematical language, or quantified module parameters for the following statements.

\paragraph{Context 5 (line 5)}
\noindent\textbf{EasyCrypt name.}
\begin{lstlisting}[language=EasyCrypt,basicstyle=\ttfamily\scriptsize]
imported theories
\end{lstlisting}
\noindent\textbf{Checked EasyCrypt statement.}
\begin{lstlisting}[language=EasyCrypt,basicstyle=\ttfamily\scriptsize]
require import HAETAE_Transcript HAETAE_ROM_Programming.
\end{lstlisting}
\noindent\textbf{Formal mathematical reading.}
Mathematical context. This line imports or specializes earlier theories, making their carrier sets, functions, modules, and theorems available to the current file.

\noindent\textbf{Role in the proof.}
Use in this chapter. It fixes the proof context, imported mathematical language, or quantified module parameters for the following statements.

\paragraph{Context 6 (line 6)}
\noindent\textbf{EasyCrypt name.}
\begin{lstlisting}[language=EasyCrypt,basicstyle=\ttfamily\scriptsize]
imported theories
\end{lstlisting}
\noindent\textbf{Checked EasyCrypt statement.}
\begin{lstlisting}[language=EasyCrypt,basicstyle=\ttfamily\scriptsize]
require import HAETAE_HopGames.
\end{lstlisting}
\noindent\textbf{Formal mathematical reading.}
Mathematical context. This line imports or specializes earlier theories, making their carrier sets, functions, modules, and theorems available to the current file.

\noindent\textbf{Role in the proof.}
Use in this chapter. It fixes the proof context, imported mathematical language, or quantified module parameters for the following statements.

\paragraph{Context 7 (line 8)}
\noindent\textbf{EasyCrypt name.}
\begin{lstlisting}[language=EasyCrypt,basicstyle=\ttfamily\scriptsize]
HAETAE_Security
\end{lstlisting}
\noindent\textbf{Checked EasyCrypt statement.}
\begin{lstlisting}[language=EasyCrypt,basicstyle=\ttfamily\scriptsize]
theory HAETAE_Security.
\end{lstlisting}
\noindent\textbf{Formal mathematical reading.}
Mathematical context. This begins a namespace for the formal development in this file.

\noindent\textbf{Role in the proof.}
Use in this chapter. It fixes the proof context, imported mathematical language, or quantified module parameters for the following statements.

\paragraph{Context 8 (line 10)}
\noindent\textbf{EasyCrypt name.}
\begin{lstlisting}[language=EasyCrypt,basicstyle=\ttfamily\scriptsize]
opened namespace
\end{lstlisting}
\noindent\textbf{Checked EasyCrypt statement.}
\begin{lstlisting}[language=EasyCrypt,basicstyle=\ttfamily\scriptsize]
import HAETAE_Algebra.
\end{lstlisting}
\noindent\textbf{Formal mathematical reading.}
Mathematical context. This line imports or specializes earlier theories, making their carrier sets, functions, modules, and theorems available to the current file.

\noindent\textbf{Role in the proof.}
Use in this chapter. It fixes the proof context, imported mathematical language, or quantified module parameters for the following statements.

\paragraph{Context 9 (line 11)}
\noindent\textbf{EasyCrypt name.}
\begin{lstlisting}[language=EasyCrypt,basicstyle=\ttfamily\scriptsize]
opened namespace
\end{lstlisting}
\noindent\textbf{Checked EasyCrypt statement.}
\begin{lstlisting}[language=EasyCrypt,basicstyle=\ttfamily\scriptsize]
import HAETAE_Distributions.
\end{lstlisting}
\noindent\textbf{Formal mathematical reading.}
Mathematical context. This line imports or specializes earlier theories, making their carrier sets, functions, modules, and theorems available to the current file.

\noindent\textbf{Role in the proof.}
Use in this chapter. It fixes the proof context, imported mathematical language, or quantified module parameters for the following statements.

\paragraph{Context 10 (line 12)}
\noindent\textbf{EasyCrypt name.}
\begin{lstlisting}[language=EasyCrypt,basicstyle=\ttfamily\scriptsize]
opened namespace
\end{lstlisting}
\noindent\textbf{Checked EasyCrypt statement.}
\begin{lstlisting}[language=EasyCrypt,basicstyle=\ttfamily\scriptsize]
import HAETAE_Reductions.
\end{lstlisting}
\noindent\textbf{Formal mathematical reading.}
Mathematical context. This line imports or specializes earlier theories, making their carrier sets, functions, modules, and theorems available to the current file.

\noindent\textbf{Role in the proof.}
Use in this chapter. It fixes the proof context, imported mathematical language, or quantified module parameters for the following statements.

\paragraph{Context 11 (line 13)}
\noindent\textbf{EasyCrypt name.}
\begin{lstlisting}[language=EasyCrypt,basicstyle=\ttfamily\scriptsize]
opened namespace
\end{lstlisting}
\noindent\textbf{Checked EasyCrypt statement.}
\begin{lstlisting}[language=EasyCrypt,basicstyle=\ttfamily\scriptsize]
import HAETAE_Assumptions.
\end{lstlisting}
\noindent\textbf{Formal mathematical reading.}
Mathematical context. This line imports or specializes earlier theories, making their carrier sets, functions, modules, and theorems available to the current file.

\noindent\textbf{Role in the proof.}
Use in this chapter. It fixes the proof context, imported mathematical language, or quantified module parameters for the following statements.

\paragraph{Context 12 (line 14)}
\noindent\textbf{EasyCrypt name.}
\begin{lstlisting}[language=EasyCrypt,basicstyle=\ttfamily\scriptsize]
opened namespace
\end{lstlisting}
\noindent\textbf{Checked EasyCrypt statement.}
\begin{lstlisting}[language=EasyCrypt,basicstyle=\ttfamily\scriptsize]
import HAETAE_Rejection.
\end{lstlisting}
\noindent\textbf{Formal mathematical reading.}
Mathematical context. This line imports or specializes earlier theories, making their carrier sets, functions, modules, and theorems available to the current file.

\noindent\textbf{Role in the proof.}
Use in this chapter. It fixes the proof context, imported mathematical language, or quantified module parameters for the following statements.

\paragraph{Context 13 (line 15)}
\noindent\textbf{EasyCrypt name.}
\begin{lstlisting}[language=EasyCrypt,basicstyle=\ttfamily\scriptsize]
opened namespace
\end{lstlisting}
\noindent\textbf{Checked EasyCrypt statement.}
\begin{lstlisting}[language=EasyCrypt,basicstyle=\ttfamily\scriptsize]
import HAETAE_Transcript.
\end{lstlisting}
\noindent\textbf{Formal mathematical reading.}
Mathematical context. This line imports or specializes earlier theories, making their carrier sets, functions, modules, and theorems available to the current file.

\noindent\textbf{Role in the proof.}
Use in this chapter. It fixes the proof context, imported mathematical language, or quantified module parameters for the following statements.

\paragraph{Context 14 (line 16)}
\noindent\textbf{EasyCrypt name.}
\begin{lstlisting}[language=EasyCrypt,basicstyle=\ttfamily\scriptsize]
opened namespace
\end{lstlisting}
\noindent\textbf{Checked EasyCrypt statement.}
\begin{lstlisting}[language=EasyCrypt,basicstyle=\ttfamily\scriptsize]
import HAETAE_ROM_Programming.
\end{lstlisting}
\noindent\textbf{Formal mathematical reading.}
Mathematical context. This line imports or specializes earlier theories, making their carrier sets, functions, modules, and theorems available to the current file.

\noindent\textbf{Role in the proof.}
Use in this chapter. It fixes the proof context, imported mathematical language, or quantified module parameters for the following statements.

\paragraph{Context 15 (line 17)}
\noindent\textbf{EasyCrypt name.}
\begin{lstlisting}[language=EasyCrypt,basicstyle=\ttfamily\scriptsize]
opened namespace
\end{lstlisting}
\noindent\textbf{Checked EasyCrypt statement.}
\begin{lstlisting}[language=EasyCrypt,basicstyle=\ttfamily\scriptsize]
import HAETAE_HopGames.
\end{lstlisting}
\noindent\textbf{Formal mathematical reading.}
Mathematical context. This line imports or specializes earlier theories, making their carrier sets, functions, modules, and theorems available to the current file.

\noindent\textbf{Role in the proof.}
Use in this chapter. It fixes the proof context, imported mathematical language, or quantified module parameters for the following statements.

\paragraph{Context 16 (line 18)}
\noindent\textbf{EasyCrypt name.}
\begin{lstlisting}[language=EasyCrypt,basicstyle=\ttfamily\scriptsize]
opened namespace
\end{lstlisting}
\noindent\textbf{Checked EasyCrypt statement.}
\begin{lstlisting}[language=EasyCrypt,basicstyle=\ttfamily\scriptsize]
import HAETAE_Scheme.
\end{lstlisting}
\noindent\textbf{Formal mathematical reading.}
Mathematical context. This line imports or specializes earlier theories, making their carrier sets, functions, modules, and theorems available to the current file.

\noindent\textbf{Role in the proof.}
Use in this chapter. It fixes the proof context, imported mathematical language, or quantified module parameters for the following statements.

\paragraph{Context 17 (line 19)}
\noindent\textbf{EasyCrypt name.}
\begin{lstlisting}[language=EasyCrypt,basicstyle=\ttfamily\scriptsize]
opened namespace
\end{lstlisting}
\noindent\textbf{Checked EasyCrypt statement.}
\begin{lstlisting}[language=EasyCrypt,basicstyle=\ttfamily\scriptsize]
import RealOrder.
\end{lstlisting}
\noindent\textbf{Formal mathematical reading.}
Mathematical context. This line imports or specializes earlier theories, making their carrier sets, functions, modules, and theorems available to the current file.

\noindent\textbf{Role in the proof.}
Use in this chapter. It fixes the proof context, imported mathematical language, or quantified module parameters for the following statements.

\paragraph{Context 18 (line 35)}
\noindent\textbf{EasyCrypt name.}
\begin{lstlisting}[language=EasyCrypt,basicstyle=\ttfamily\scriptsize]
section_line_35
\end{lstlisting}
\noindent\textbf{Checked EasyCrypt statement.}
\begin{lstlisting}[language=EasyCrypt,basicstyle=\ttfamily\scriptsize]
section.
\end{lstlisting}
\noindent\textbf{Formal mathematical reading.}
Mathematical context. This opens a scoped block of hypotheses, module parameters, and lemmas.

\noindent\textbf{Role in the proof.}
Use in this chapter. It fixes the proof context, imported mathematical language, or quantified module parameters for the following statements.

\paragraph{Context 19 (line 37)}
\noindent\textbf{EasyCrypt name.}
\begin{lstlisting}[language=EasyCrypt,basicstyle=\ttfamily\scriptsize]
H
\end{lstlisting}
\noindent\textbf{Checked EasyCrypt statement.}
\begin{lstlisting}[language=EasyCrypt,basicstyle=\ttfamily\scriptsize]
declare module H <: SIG.Oracle.
\end{lstlisting}
\noindent\textbf{Formal mathematical reading.}
Mathematical context. This is universal quantification over an adversary, oracle, scheme, sampler, or reduction module satisfying the stated interface and memory restrictions.

\noindent\textbf{Role in the proof.}
Use in this chapter. It fixes the proof context, imported mathematical language, or quantified module parameters for the following statements.

\paragraph{Context 20 (line 38)}
\noindent\textbf{EasyCrypt name.}
\begin{lstlisting}[language=EasyCrypt,basicstyle=\ttfamily\scriptsize]
A
\end{lstlisting}
\noindent\textbf{Checked EasyCrypt statement.}
\begin{lstlisting}[language=EasyCrypt,basicstyle=\ttfamily\scriptsize]
declare module A <: SIG.CORR_ADV {-H, -HAETAE}.
\end{lstlisting}
\noindent\textbf{Formal mathematical reading.}
Mathematical context. This is universal quantification over an adversary, oracle, scheme, sampler, or reduction module satisfying the stated interface and memory restrictions.

\noindent\textbf{Role in the proof.}
Use in this chapter. It fixes the proof context, imported mathematical language, or quantified module parameters for the following statements.

\paragraph{Context 21 (line 73)}
\noindent\textbf{EasyCrypt name.}
\begin{lstlisting}[language=EasyCrypt,basicstyle=\ttfamily\scriptsize]
NoSigningSecurity
\end{lstlisting}
\noindent\textbf{Checked EasyCrypt statement.}
\begin{lstlisting}[language=EasyCrypt,basicstyle=\ttfamily\scriptsize]
section NoSigningSecurity.
\end{lstlisting}
\noindent\textbf{Formal mathematical reading.}
Mathematical context. This opens a scoped block of hypotheses, module parameters, and lemmas.

\noindent\textbf{Role in the proof.}
Use in this chapter. It fixes the proof context, imported mathematical language, or quantified module parameters for the following statements.

\paragraph{Context 22 (line 75)}
\noindent\textbf{EasyCrypt name.}
\begin{lstlisting}[language=EasyCrypt,basicstyle=\ttfamily\scriptsize]
H
\end{lstlisting}
\noindent\textbf{Checked EasyCrypt statement.}
\begin{lstlisting}[language=EasyCrypt,basicstyle=\ttfamily\scriptsize]
declare module H <: SIG.Oracle {-SIG.EUF_CMA}.
\end{lstlisting}
\noindent\textbf{Formal mathematical reading.}
Mathematical context. This is universal quantification over an adversary, oracle, scheme, sampler, or reduction module satisfying the stated interface and memory restrictions.

\noindent\textbf{Role in the proof.}
Use in this chapter. It fixes the proof context, imported mathematical language, or quantified module parameters for the following statements.

\paragraph{Context 23 (line 76)}
\noindent\textbf{EasyCrypt name.}
\begin{lstlisting}[language=EasyCrypt,basicstyle=\ttfamily\scriptsize]
N
\end{lstlisting}
\noindent\textbf{Checked EasyCrypt statement.}
\begin{lstlisting}[language=EasyCrypt,basicstyle=\ttfamily\scriptsize]
declare module N <: SIG.NMA_Adversary {-H, -HAETAE, -SIG.EUF_CMA}.
\end{lstlisting}
\noindent\textbf{Formal mathematical reading.}
Mathematical context. This is universal quantification over an adversary, oracle, scheme, sampler, or reduction module satisfying the stated interface and memory restrictions.

\noindent\textbf{Role in the proof.}
Use in this chapter. It fixes the proof context, imported mathematical language, or quantified module parameters for the following statements.

\paragraph{Context 24 (line 77)}
\noindent\textbf{EasyCrypt name.}
\begin{lstlisting}[language=EasyCrypt,basicstyle=\ttfamily\scriptsize]
Bmlwe
\end{lstlisting}
\noindent\textbf{Checked EasyCrypt statement.}
\begin{lstlisting}[language=EasyCrypt,basicstyle=\ttfamily\scriptsize]
declare module Bmlwe <: MLWE_Reduction.
\end{lstlisting}
\noindent\textbf{Formal mathematical reading.}
Mathematical context. This is universal quantification over an adversary, oracle, scheme, sampler, or reduction module satisfying the stated interface and memory restrictions.

\noindent\textbf{Role in the proof.}
Use in this chapter. It fixes the proof context, imported mathematical language, or quantified module parameters for the following statements.

\paragraph{Context 25 (line 78)}
\noindent\textbf{EasyCrypt name.}
\begin{lstlisting}[language=EasyCrypt,basicstyle=\ttfamily\scriptsize]
Bbimodal
\end{lstlisting}
\noindent\textbf{Checked EasyCrypt statement.}
\begin{lstlisting}[language=EasyCrypt,basicstyle=\ttfamily\scriptsize]
declare module Bbimodal <: BimodalSelfTargetMSIS_Reduction.
\end{lstlisting}
\noindent\textbf{Formal mathematical reading.}
Mathematical context. This is universal quantification over an adversary, oracle, scheme, sampler, or reduction module satisfying the stated interface and memory restrictions.

\noindent\textbf{Role in the proof.}
Use in this chapter. It fixes the proof context, imported mathematical language, or quantified module parameters for the following statements.

\paragraph{Context 26 (line 79)}
\noindent\textbf{EasyCrypt name.}
\begin{lstlisting}[language=EasyCrypt,basicstyle=\ttfamily\scriptsize]
Bsis
\end{lstlisting}
\noindent\textbf{Checked EasyCrypt statement.}
\begin{lstlisting}[language=EasyCrypt,basicstyle=\ttfamily\scriptsize]
declare module Bsis <: ModuleSIS_Reduction.
\end{lstlisting}
\noindent\textbf{Formal mathematical reading.}
Mathematical context. This is universal quantification over an adversary, oracle, scheme, sampler, or reduction module satisfying the stated interface and memory restrictions.

\noindent\textbf{Role in the proof.}
Use in this chapter. It fixes the proof context, imported mathematical language, or quantified module parameters for the following statements.

\paragraph{Context 27 (line 127)}
\noindent\textbf{EasyCrypt name.}
\begin{lstlisting}[language=EasyCrypt,basicstyle=\ttfamily\scriptsize]
section_line_127
\end{lstlisting}
\noindent\textbf{Checked EasyCrypt statement.}
\begin{lstlisting}[language=EasyCrypt,basicstyle=\ttfamily\scriptsize]
section.
\end{lstlisting}
\noindent\textbf{Formal mathematical reading.}
Mathematical context. This opens a scoped block of hypotheses, module parameters, and lemmas.

\noindent\textbf{Role in the proof.}
Use in this chapter. It fixes the proof context, imported mathematical language, or quantified module parameters for the following statements.

\paragraph{Context 28 (line 129)}
\noindent\textbf{EasyCrypt name.}
\begin{lstlisting}[language=EasyCrypt,basicstyle=\ttfamily\scriptsize]
H
\end{lstlisting}
\noindent\textbf{Checked EasyCrypt statement.}
\begin{lstlisting}[language=EasyCrypt,basicstyle=\ttfamily\scriptsize]
declare module H <: SIG.Oracle.
\end{lstlisting}
\noindent\textbf{Formal mathematical reading.}
Mathematical context. This is universal quantification over an adversary, oracle, scheme, sampler, or reduction module satisfying the stated interface and memory restrictions.

\noindent\textbf{Role in the proof.}
Use in this chapter. It fixes the proof context, imported mathematical language, or quantified module parameters for the following statements.

\paragraph{Context 29 (line 130)}
\noindent\textbf{EasyCrypt name.}
\begin{lstlisting}[language=EasyCrypt,basicstyle=\ttfamily\scriptsize]
A
\end{lstlisting}
\noindent\textbf{Checked EasyCrypt statement.}
\begin{lstlisting}[language=EasyCrypt,basicstyle=\ttfamily\scriptsize]
declare module A <: SIG.Adversary {-H, -HAETAE}.
\end{lstlisting}
\noindent\textbf{Formal mathematical reading.}
Mathematical context. This is universal quantification over an adversary, oracle, scheme, sampler, or reduction module satisfying the stated interface and memory restrictions.

\noindent\textbf{Role in the proof.}
Use in this chapter. It fixes the proof context, imported mathematical language, or quantified module parameters for the following statements.

\paragraph{Context 30 (line 131)}
\noindent\textbf{EasyCrypt name.}
\begin{lstlisting}[language=EasyCrypt,basicstyle=\ttfamily\scriptsize]
N
\end{lstlisting}
\noindent\textbf{Checked EasyCrypt statement.}
\begin{lstlisting}[language=EasyCrypt,basicstyle=\ttfamily\scriptsize]
declare module N <: SIG.NMA_Adversary {-H, -HAETAE}.
\end{lstlisting}
\noindent\textbf{Formal mathematical reading.}
Mathematical context. This is universal quantification over an adversary, oracle, scheme, sampler, or reduction module satisfying the stated interface and memory restrictions.

\noindent\textbf{Role in the proof.}
Use in this chapter. It fixes the proof context, imported mathematical language, or quantified module parameters for the following statements.

\paragraph{Context 31 (line 132)}
\noindent\textbf{EasyCrypt name.}
\begin{lstlisting}[language=EasyCrypt,basicstyle=\ttfamily\scriptsize]
Bmlwe
\end{lstlisting}
\noindent\textbf{Checked EasyCrypt statement.}
\begin{lstlisting}[language=EasyCrypt,basicstyle=\ttfamily\scriptsize]
declare module Bmlwe <: MLWE_Reduction.
\end{lstlisting}
\noindent\textbf{Formal mathematical reading.}
Mathematical context. This is universal quantification over an adversary, oracle, scheme, sampler, or reduction module satisfying the stated interface and memory restrictions.

\noindent\textbf{Role in the proof.}
Use in this chapter. It fixes the proof context, imported mathematical language, or quantified module parameters for the following statements.

\paragraph{Context 32 (line 133)}
\noindent\textbf{EasyCrypt name.}
\begin{lstlisting}[language=EasyCrypt,basicstyle=\ttfamily\scriptsize]
Bbimodal
\end{lstlisting}
\noindent\textbf{Checked EasyCrypt statement.}
\begin{lstlisting}[language=EasyCrypt,basicstyle=\ttfamily\scriptsize]
declare module Bbimodal <: BimodalSelfTargetMSIS_Reduction.
\end{lstlisting}
\noindent\textbf{Formal mathematical reading.}
Mathematical context. This is universal quantification over an adversary, oracle, scheme, sampler, or reduction module satisfying the stated interface and memory restrictions.

\noindent\textbf{Role in the proof.}
Use in this chapter. It fixes the proof context, imported mathematical language, or quantified module parameters for the following statements.

\paragraph{Context 33 (line 134)}
\noindent\textbf{EasyCrypt name.}
\begin{lstlisting}[language=EasyCrypt,basicstyle=\ttfamily\scriptsize]
Bsis
\end{lstlisting}
\noindent\textbf{Checked EasyCrypt statement.}
\begin{lstlisting}[language=EasyCrypt,basicstyle=\ttfamily\scriptsize]
declare module Bsis <: ModuleSIS_Reduction.
\end{lstlisting}
\noindent\textbf{Formal mathematical reading.}
Mathematical context. This is universal quantification over an adversary, oracle, scheme, sampler, or reduction module satisfying the stated interface and memory restrictions.

\noindent\textbf{Role in the proof.}
Use in this chapter. It fixes the proof context, imported mathematical language, or quantified module parameters for the following statements.

\paragraph{Context 34 (line 211)}
\noindent\textbf{EasyCrypt name.}
\begin{lstlisting}[language=EasyCrypt,basicstyle=\ttfamily\scriptsize]
FullReductionWithSimulatorHop
\end{lstlisting}
\noindent\textbf{Checked EasyCrypt statement.}
\begin{lstlisting}[language=EasyCrypt,basicstyle=\ttfamily\scriptsize]
section FullReductionWithSimulatorHop.
\end{lstlisting}
\noindent\textbf{Formal mathematical reading.}
Mathematical context. This opens a scoped block of hypotheses, module parameters, and lemmas.

\noindent\textbf{Role in the proof.}
Use in this chapter. It fixes the proof context, imported mathematical language, or quantified module parameters for the following statements.

\paragraph{Context 35 (line 213)}
\noindent\textbf{EasyCrypt name.}
\begin{lstlisting}[language=EasyCrypt,basicstyle=\ttfamily\scriptsize]
H
\end{lstlisting}
\noindent\textbf{Checked EasyCrypt statement.}
\begin{lstlisting}[language=EasyCrypt,basicstyle=\ttfamily\scriptsize]
declare module H <: SIG.Oracle {-SIG.EUF_CMA, -EUF_CMA_SimulatedSign,
                                 -RealSigningSimulator}.
\end{lstlisting}
\noindent\textbf{Formal mathematical reading.}
Mathematical context. This is universal quantification over an adversary, oracle, scheme, sampler, or reduction module satisfying the stated interface and memory restrictions.

\noindent\textbf{Role in the proof.}
Use in this chapter. It fixes the proof context, imported mathematical language, or quantified module parameters for the following statements.

\paragraph{Context 36 (line 215)}
\noindent\textbf{EasyCrypt name.}
\begin{lstlisting}[language=EasyCrypt,basicstyle=\ttfamily\scriptsize]
A
\end{lstlisting}
\noindent\textbf{Checked EasyCrypt statement.}
\begin{lstlisting}[language=EasyCrypt,basicstyle=\ttfamily\scriptsize]
declare module A <: SIG.Adversary {-H, -HAETAE, -SIG.EUF_CMA,
                                    -EUF_CMA_SimulatedSign,
                                    -RealSigningSimulator}.
\end{lstlisting}
\noindent\textbf{Formal mathematical reading.}
Mathematical context. This is universal quantification over an adversary, oracle, scheme, sampler, or reduction module satisfying the stated interface and memory restrictions.

\noindent\textbf{Role in the proof.}
Use in this chapter. It fixes the proof context, imported mathematical language, or quantified module parameters for the following statements.

\paragraph{Context 37 (line 218)}
\noindent\textbf{EasyCrypt name.}
\begin{lstlisting}[language=EasyCrypt,basicstyle=\ttfamily\scriptsize]
N
\end{lstlisting}
\noindent\textbf{Checked EasyCrypt statement.}
\begin{lstlisting}[language=EasyCrypt,basicstyle=\ttfamily\scriptsize]
declare module N <: SIG.NMA_Adversary {-H, -HAETAE}.
\end{lstlisting}
\noindent\textbf{Formal mathematical reading.}
Mathematical context. This is universal quantification over an adversary, oracle, scheme, sampler, or reduction module satisfying the stated interface and memory restrictions.

\noindent\textbf{Role in the proof.}
Use in this chapter. It fixes the proof context, imported mathematical language, or quantified module parameters for the following statements.

\paragraph{Context 38 (line 219)}
\noindent\textbf{EasyCrypt name.}
\begin{lstlisting}[language=EasyCrypt,basicstyle=\ttfamily\scriptsize]
Bmlwe
\end{lstlisting}
\noindent\textbf{Checked EasyCrypt statement.}
\begin{lstlisting}[language=EasyCrypt,basicstyle=\ttfamily\scriptsize]
declare module Bmlwe <: MLWE_Reduction.
\end{lstlisting}
\noindent\textbf{Formal mathematical reading.}
Mathematical context. This is universal quantification over an adversary, oracle, scheme, sampler, or reduction module satisfying the stated interface and memory restrictions.

\noindent\textbf{Role in the proof.}
Use in this chapter. It fixes the proof context, imported mathematical language, or quantified module parameters for the following statements.

\paragraph{Context 39 (line 220)}
\noindent\textbf{EasyCrypt name.}
\begin{lstlisting}[language=EasyCrypt,basicstyle=\ttfamily\scriptsize]
Bbimodal
\end{lstlisting}
\noindent\textbf{Checked EasyCrypt statement.}
\begin{lstlisting}[language=EasyCrypt,basicstyle=\ttfamily\scriptsize]
declare module Bbimodal <: BimodalSelfTargetMSIS_Reduction.
\end{lstlisting}
\noindent\textbf{Formal mathematical reading.}
Mathematical context. This is universal quantification over an adversary, oracle, scheme, sampler, or reduction module satisfying the stated interface and memory restrictions.

\noindent\textbf{Role in the proof.}
Use in this chapter. It fixes the proof context, imported mathematical language, or quantified module parameters for the following statements.

\paragraph{Context 40 (line 221)}
\noindent\textbf{EasyCrypt name.}
\begin{lstlisting}[language=EasyCrypt,basicstyle=\ttfamily\scriptsize]
Bsis
\end{lstlisting}
\noindent\textbf{Checked EasyCrypt statement.}
\begin{lstlisting}[language=EasyCrypt,basicstyle=\ttfamily\scriptsize]
declare module Bsis <: ModuleSIS_Reduction.
\end{lstlisting}
\noindent\textbf{Formal mathematical reading.}
Mathematical context. This is universal quantification over an adversary, oracle, scheme, sampler, or reduction module satisfying the stated interface and memory restrictions.

\noindent\textbf{Role in the proof.}
Use in this chapter. It fixes the proof context, imported mathematical language, or quantified module parameters for the following statements.

\paragraph{Context 41 (line 306)}
\noindent\textbf{EasyCrypt name.}
\begin{lstlisting}[language=EasyCrypt,basicstyle=\ttfamily\scriptsize]
FullReductionWithROMTranscriptHop
\end{lstlisting}
\noindent\textbf{Checked EasyCrypt statement.}
\begin{lstlisting}[language=EasyCrypt,basicstyle=\ttfamily\scriptsize]
section FullReductionWithROMTranscriptHop.
\end{lstlisting}
\noindent\textbf{Formal mathematical reading.}
Mathematical context. This opens a scoped block of hypotheses, module parameters, and lemmas.

\noindent\textbf{Role in the proof.}
Use in this chapter. It fixes the proof context, imported mathematical language, or quantified module parameters for the following statements.

\paragraph{Context 42 (line 308)}
\noindent\textbf{EasyCrypt name.}
\begin{lstlisting}[language=EasyCrypt,basicstyle=\ttfamily\scriptsize]
H
\end{lstlisting}
\noindent\textbf{Checked EasyCrypt statement.}
\begin{lstlisting}[language=EasyCrypt,basicstyle=\ttfamily\scriptsize]
declare module H <: SIG.Oracle {-SIG.EUF_CMA,
                                 -EUF_CMA_InternalTranscriptSign,
                                 -EUF_CMA_ROMInternalTranscriptSign}.
\end{lstlisting}
\noindent\textbf{Formal mathematical reading.}
Mathematical context. This is universal quantification over an adversary, oracle, scheme, sampler, or reduction module satisfying the stated interface and memory restrictions.

\noindent\textbf{Role in the proof.}
Use in this chapter. It fixes the proof context, imported mathematical language, or quantified module parameters for the following statements.

\paragraph{Context 43 (line 311)}
\noindent\textbf{EasyCrypt name.}
\begin{lstlisting}[language=EasyCrypt,basicstyle=\ttfamily\scriptsize]
A
\end{lstlisting}
\noindent\textbf{Checked EasyCrypt statement.}
\begin{lstlisting}[language=EasyCrypt,basicstyle=\ttfamily\scriptsize]
declare module A <: SIG.Adversary {-H, -HAETAE, -SIG.EUF_CMA,
                                    -EUF_CMA_InternalTranscriptSign,
                                    -EUF_CMA_ROMInternalTranscriptSign}.
\end{lstlisting}
\noindent\textbf{Formal mathematical reading.}
Mathematical context. This is universal quantification over an adversary, oracle, scheme, sampler, or reduction module satisfying the stated interface and memory restrictions.

\noindent\textbf{Role in the proof.}
Use in this chapter. It fixes the proof context, imported mathematical language, or quantified module parameters for the following statements.

\paragraph{Context 44 (line 314)}
\noindent\textbf{EasyCrypt name.}
\begin{lstlisting}[language=EasyCrypt,basicstyle=\ttfamily\scriptsize]
N
\end{lstlisting}
\noindent\textbf{Checked EasyCrypt statement.}
\begin{lstlisting}[language=EasyCrypt,basicstyle=\ttfamily\scriptsize]
declare module N <: SIG.NMA_Adversary {-H, -HAETAE}.
\end{lstlisting}
\noindent\textbf{Formal mathematical reading.}
Mathematical context. This is universal quantification over an adversary, oracle, scheme, sampler, or reduction module satisfying the stated interface and memory restrictions.

\noindent\textbf{Role in the proof.}
Use in this chapter. It fixes the proof context, imported mathematical language, or quantified module parameters for the following statements.

\paragraph{Context 45 (line 315)}
\noindent\textbf{EasyCrypt name.}
\begin{lstlisting}[language=EasyCrypt,basicstyle=\ttfamily\scriptsize]
Bmlwe
\end{lstlisting}
\noindent\textbf{Checked EasyCrypt statement.}
\begin{lstlisting}[language=EasyCrypt,basicstyle=\ttfamily\scriptsize]
declare module Bmlwe <: MLWE_Reduction.
\end{lstlisting}
\noindent\textbf{Formal mathematical reading.}
Mathematical context. This is universal quantification over an adversary, oracle, scheme, sampler, or reduction module satisfying the stated interface and memory restrictions.

\noindent\textbf{Role in the proof.}
Use in this chapter. It fixes the proof context, imported mathematical language, or quantified module parameters for the following statements.

\paragraph{Context 46 (line 316)}
\noindent\textbf{EasyCrypt name.}
\begin{lstlisting}[language=EasyCrypt,basicstyle=\ttfamily\scriptsize]
Bbimodal
\end{lstlisting}
\noindent\textbf{Checked EasyCrypt statement.}
\begin{lstlisting}[language=EasyCrypt,basicstyle=\ttfamily\scriptsize]
declare module Bbimodal <: BimodalSelfTargetMSIS_Reduction.
\end{lstlisting}
\noindent\textbf{Formal mathematical reading.}
Mathematical context. This is universal quantification over an adversary, oracle, scheme, sampler, or reduction module satisfying the stated interface and memory restrictions.

\noindent\textbf{Role in the proof.}
Use in this chapter. It fixes the proof context, imported mathematical language, or quantified module parameters for the following statements.

\paragraph{Context 47 (line 317)}
\noindent\textbf{EasyCrypt name.}
\begin{lstlisting}[language=EasyCrypt,basicstyle=\ttfamily\scriptsize]
Bsis
\end{lstlisting}
\noindent\textbf{Checked EasyCrypt statement.}
\begin{lstlisting}[language=EasyCrypt,basicstyle=\ttfamily\scriptsize]
declare module Bsis <: ModuleSIS_Reduction.
\end{lstlisting}
\noindent\textbf{Formal mathematical reading.}
Mathematical context. This is universal quantification over an adversary, oracle, scheme, sampler, or reduction module satisfying the stated interface and memory restrictions.

\noindent\textbf{Role in the proof.}
Use in this chapter. It fixes the proof context, imported mathematical language, or quantified module parameters for the following statements.

\paragraph{Context 48 (line 407)}
\noindent\textbf{EasyCrypt name.}
\begin{lstlisting}[language=EasyCrypt,basicstyle=\ttfamily\scriptsize]
MechanizedBimodalModuleSISLift
\end{lstlisting}
\noindent\textbf{Checked EasyCrypt statement.}
\begin{lstlisting}[language=EasyCrypt,basicstyle=\ttfamily\scriptsize]
section MechanizedBimodalModuleSISLift.
\end{lstlisting}
\noindent\textbf{Formal mathematical reading.}
Mathematical context. This opens a scoped block of hypotheses, module parameters, and lemmas.

\noindent\textbf{Role in the proof.}
Use in this chapter. It fixes the proof context, imported mathematical language, or quantified module parameters for the following statements.

\paragraph{Context 49 (line 409)}
\noindent\textbf{EasyCrypt name.}
\begin{lstlisting}[language=EasyCrypt,basicstyle=\ttfamily\scriptsize]
H
\end{lstlisting}
\noindent\textbf{Checked EasyCrypt statement.}
\begin{lstlisting}[language=EasyCrypt,basicstyle=\ttfamily\scriptsize]
declare module H <: SIG.Oracle {-SIG.EUF_CMA, -EUF_CMA_SimulatedSign,
                                 -RealSigningSimulator}.
\end{lstlisting}
\noindent\textbf{Formal mathematical reading.}
Mathematical context. This is universal quantification over an adversary, oracle, scheme, sampler, or reduction module satisfying the stated interface and memory restrictions.

\noindent\textbf{Role in the proof.}
Use in this chapter. It fixes the proof context, imported mathematical language, or quantified module parameters for the following statements.

\paragraph{Context 50 (line 411)}
\noindent\textbf{EasyCrypt name.}
\begin{lstlisting}[language=EasyCrypt,basicstyle=\ttfamily\scriptsize]
A
\end{lstlisting}
\noindent\textbf{Checked EasyCrypt statement.}
\begin{lstlisting}[language=EasyCrypt,basicstyle=\ttfamily\scriptsize]
declare module A <: SIG.Adversary {-H, -HAETAE, -SIG.EUF_CMA,
                                    -EUF_CMA_SimulatedSign,
                                    -RealSigningSimulator}.
\end{lstlisting}
\noindent\textbf{Formal mathematical reading.}
Mathematical context. This is universal quantification over an adversary, oracle, scheme, sampler, or reduction module satisfying the stated interface and memory restrictions.

\noindent\textbf{Role in the proof.}
Use in this chapter. It fixes the proof context, imported mathematical language, or quantified module parameters for the following statements.

\paragraph{Context 51 (line 414)}
\noindent\textbf{EasyCrypt name.}
\begin{lstlisting}[language=EasyCrypt,basicstyle=\ttfamily\scriptsize]
N
\end{lstlisting}
\noindent\textbf{Checked EasyCrypt statement.}
\begin{lstlisting}[language=EasyCrypt,basicstyle=\ttfamily\scriptsize]
declare module N <: SIG.NMA_Adversary {-H, -HAETAE}.
\end{lstlisting}
\noindent\textbf{Formal mathematical reading.}
Mathematical context. This is universal quantification over an adversary, oracle, scheme, sampler, or reduction module satisfying the stated interface and memory restrictions.

\noindent\textbf{Role in the proof.}
Use in this chapter. It fixes the proof context, imported mathematical language, or quantified module parameters for the following statements.

\paragraph{Context 52 (line 415)}
\noindent\textbf{EasyCrypt name.}
\begin{lstlisting}[language=EasyCrypt,basicstyle=\ttfamily\scriptsize]
Bmlwe
\end{lstlisting}
\noindent\textbf{Checked EasyCrypt statement.}
\begin{lstlisting}[language=EasyCrypt,basicstyle=\ttfamily\scriptsize]
declare module Bmlwe <: MLWE_Reduction.
\end{lstlisting}
\noindent\textbf{Formal mathematical reading.}
Mathematical context. This is universal quantification over an adversary, oracle, scheme, sampler, or reduction module satisfying the stated interface and memory restrictions.

\noindent\textbf{Role in the proof.}
Use in this chapter. It fixes the proof context, imported mathematical language, or quantified module parameters for the following statements.

\paragraph{Context 53 (line 416)}
\noindent\textbf{EasyCrypt name.}
\begin{lstlisting}[language=EasyCrypt,basicstyle=\ttfamily\scriptsize]
Bbimodal
\end{lstlisting}
\noindent\textbf{Checked EasyCrypt statement.}
\begin{lstlisting}[language=EasyCrypt,basicstyle=\ttfamily\scriptsize]
declare module Bbimodal <: BimodalSelfTargetMSIS_Reduction.
\end{lstlisting}
\noindent\textbf{Formal mathematical reading.}
Mathematical context. This is universal quantification over an adversary, oracle, scheme, sampler, or reduction module satisfying the stated interface and memory restrictions.

\noindent\textbf{Role in the proof.}
Use in this chapter. It fixes the proof context, imported mathematical language, or quantified module parameters for the following statements.

\paragraph{Context 54 (line 461)}
\noindent\textbf{EasyCrypt name.}
\begin{lstlisting}[language=EasyCrypt,basicstyle=\ttfamily\scriptsize]
MechanizedROMTranscriptBimodalLift
\end{lstlisting}
\noindent\textbf{Checked EasyCrypt statement.}
\begin{lstlisting}[language=EasyCrypt,basicstyle=\ttfamily\scriptsize]
section MechanizedROMTranscriptBimodalLift.
\end{lstlisting}
\noindent\textbf{Formal mathematical reading.}
Mathematical context. This opens a scoped block of hypotheses, module parameters, and lemmas.

\noindent\textbf{Role in the proof.}
Use in this chapter. It fixes the proof context, imported mathematical language, or quantified module parameters for the following statements.

\paragraph{Context 55 (line 463)}
\noindent\textbf{EasyCrypt name.}
\begin{lstlisting}[language=EasyCrypt,basicstyle=\ttfamily\scriptsize]
H
\end{lstlisting}
\noindent\textbf{Checked EasyCrypt statement.}
\begin{lstlisting}[language=EasyCrypt,basicstyle=\ttfamily\scriptsize]
declare module H <: SIG.Oracle {-SIG.EUF_CMA,
                                 -EUF_CMA_InternalTranscriptSign,
                                 -EUF_CMA_ROMInternalTranscriptSign}.
\end{lstlisting}
\noindent\textbf{Formal mathematical reading.}
Mathematical context. This is universal quantification over an adversary, oracle, scheme, sampler, or reduction module satisfying the stated interface and memory restrictions.

\noindent\textbf{Role in the proof.}
Use in this chapter. It fixes the proof context, imported mathematical language, or quantified module parameters for the following statements.

\paragraph{Context 56 (line 466)}
\noindent\textbf{EasyCrypt name.}
\begin{lstlisting}[language=EasyCrypt,basicstyle=\ttfamily\scriptsize]
A
\end{lstlisting}
\noindent\textbf{Checked EasyCrypt statement.}
\begin{lstlisting}[language=EasyCrypt,basicstyle=\ttfamily\scriptsize]
declare module A <: SIG.Adversary {-H, -HAETAE, -SIG.EUF_CMA,
                                    -EUF_CMA_InternalTranscriptSign,
                                    -EUF_CMA_ROMInternalTranscriptSign}.
\end{lstlisting}
\noindent\textbf{Formal mathematical reading.}
Mathematical context. This is universal quantification over an adversary, oracle, scheme, sampler, or reduction module satisfying the stated interface and memory restrictions.

\noindent\textbf{Role in the proof.}
Use in this chapter. It fixes the proof context, imported mathematical language, or quantified module parameters for the following statements.

\paragraph{Context 57 (line 469)}
\noindent\textbf{EasyCrypt name.}
\begin{lstlisting}[language=EasyCrypt,basicstyle=\ttfamily\scriptsize]
N
\end{lstlisting}
\noindent\textbf{Checked EasyCrypt statement.}
\begin{lstlisting}[language=EasyCrypt,basicstyle=\ttfamily\scriptsize]
declare module N <: SIG.NMA_Adversary {-H, -HAETAE}.
\end{lstlisting}
\noindent\textbf{Formal mathematical reading.}
Mathematical context. This is universal quantification over an adversary, oracle, scheme, sampler, or reduction module satisfying the stated interface and memory restrictions.

\noindent\textbf{Role in the proof.}
Use in this chapter. It fixes the proof context, imported mathematical language, or quantified module parameters for the following statements.

\paragraph{Context 58 (line 470)}
\noindent\textbf{EasyCrypt name.}
\begin{lstlisting}[language=EasyCrypt,basicstyle=\ttfamily\scriptsize]
Bmlwe
\end{lstlisting}
\noindent\textbf{Checked EasyCrypt statement.}
\begin{lstlisting}[language=EasyCrypt,basicstyle=\ttfamily\scriptsize]
declare module Bmlwe <: MLWE_Reduction.
\end{lstlisting}
\noindent\textbf{Formal mathematical reading.}
Mathematical context. This is universal quantification over an adversary, oracle, scheme, sampler, or reduction module satisfying the stated interface and memory restrictions.

\noindent\textbf{Role in the proof.}
Use in this chapter. It fixes the proof context, imported mathematical language, or quantified module parameters for the following statements.

\paragraph{Context 59 (line 471)}
\noindent\textbf{EasyCrypt name.}
\begin{lstlisting}[language=EasyCrypt,basicstyle=\ttfamily\scriptsize]
Bbimodal
\end{lstlisting}
\noindent\textbf{Checked EasyCrypt statement.}
\begin{lstlisting}[language=EasyCrypt,basicstyle=\ttfamily\scriptsize]
declare module Bbimodal <: BimodalSelfTargetMSIS_Reduction.
\end{lstlisting}
\noindent\textbf{Formal mathematical reading.}
Mathematical context. This is universal quantification over an adversary, oracle, scheme, sampler, or reduction module satisfying the stated interface and memory restrictions.

\noindent\textbf{Role in the proof.}
Use in this chapter. It fixes the proof context, imported mathematical language, or quantified module parameters for the following statements.

\paragraph{Context 60 (line 598)}
\noindent\textbf{EasyCrypt name.}
\begin{lstlisting}[language=EasyCrypt,basicstyle=\ttfamily\scriptsize]
MechanizedROMTranscriptConstructedNMALift
\end{lstlisting}
\noindent\textbf{Checked EasyCrypt statement.}
\begin{lstlisting}[language=EasyCrypt,basicstyle=\ttfamily\scriptsize]
section MechanizedROMTranscriptConstructedNMALift.
\end{lstlisting}
\noindent\textbf{Formal mathematical reading.}
Mathematical context. This opens a scoped block of hypotheses, module parameters, and lemmas.

\noindent\textbf{Role in the proof.}
Use in this chapter. It fixes the proof context, imported mathematical language, or quantified module parameters for the following statements.

\paragraph{Context 61 (line 600)}
\noindent\textbf{EasyCrypt name.}
\begin{lstlisting}[language=EasyCrypt,basicstyle=\ttfamily\scriptsize]
H
\end{lstlisting}
\noindent\textbf{Checked EasyCrypt statement.}
\begin{lstlisting}[language=EasyCrypt,basicstyle=\ttfamily\scriptsize]
declare module H <: SIG.Oracle {-SIG.EUF_CMA,
                                 -EUF_CMA_InternalTranscriptSign,
                                 -EUF_CMA_ROMInternalTranscriptSign,
                                 -ROMInternalTranscriptAsNMA,
                                 -ROMInternalTranscriptPublicSimAsNMA,
                                 -ROMInternalTranscriptPaperSimAsNMA,
                                 -ROMInternalTranscriptBudgetedPaperSimAsNMA,
                                 -ROMPaperSimSigningSampler,
                                 -RealSigningPaperSimSampler,
                                 -ROSigningAttemptPaperSimSampler,
                                 -HAETAERejectionPaperSimSampler,
                                 -ExactHyperballHAETAERejectionPaperSimSampler,
                                 -ExactHyperballPaperSimSampler,
                                 -ROExactHyperballPaperSimSampler}.
\end{lstlisting}
\noindent\textbf{Formal mathematical reading.}
Mathematical context. This is universal quantification over an adversary, oracle, scheme, sampler, or reduction module satisfying the stated interface and memory restrictions.

\noindent\textbf{Role in the proof.}
Use in this chapter. It fixes the proof context, imported mathematical language, or quantified module parameters for the following statements.

\paragraph{Context 62 (line 614)}
\noindent\textbf{EasyCrypt name.}
\begin{lstlisting}[language=EasyCrypt,basicstyle=\ttfamily\scriptsize]
A
\end{lstlisting}
\noindent\textbf{Checked EasyCrypt statement.}
\begin{lstlisting}[language=EasyCrypt,basicstyle=\ttfamily\scriptsize]
declare module A <: SIG.Adversary {-H, -HAETAE, -SIG.EUF_CMA,
                                    -EUF_CMA_InternalTranscriptSign,
                                    -EUF_CMA_ROMInternalTranscriptSign,
                                    -ROMInternalTranscriptAsNMA,
                                    -ROMInternalTranscriptPublicSimAsNMA,
                                    -ROMInternalTranscriptPaperSimAsNMA,
                                    -ROMInternalTranscriptBudgetedPaperSimAsNMA,
                                    -ROMPaperSimSigningSampler,
                                    -RealSigningPaperSimSampler,
                                    -ROSigningAttemptPaperSimSampler,
                                    -HAETAERejectionPaperSimSampler,
                                    -ExactHyperballHAETAERejectionPaperSimSampler,
                                    -ExactHyperballPaperSimSampler,
                                    -ROExactHyperballPaperSimSampler}.
\end{lstlisting}
\noindent\textbf{Formal mathematical reading.}
Mathematical context. This is universal quantification over an adversary, oracle, scheme, sampler, or reduction module satisfying the stated interface and memory restrictions.

\noindent\textbf{Role in the proof.}
Use in this chapter. It fixes the proof context, imported mathematical language, or quantified module parameters for the following statements.

\paragraph{Context 63 (line 628)}
\noindent\textbf{EasyCrypt name.}
\begin{lstlisting}[language=EasyCrypt,basicstyle=\ttfamily\scriptsize]
Samp
\end{lstlisting}
\noindent\textbf{Checked EasyCrypt statement.}
\begin{lstlisting}[language=EasyCrypt,basicstyle=\ttfamily\scriptsize]
declare module Samp <: PaperSimSigningSampler {-H, -A,
                                               -ROMInternalTranscriptPaperSimAsNMA}.
\end{lstlisting}
\noindent\textbf{Formal mathematical reading.}
Mathematical context. This is universal quantification over an adversary, oracle, scheme, sampler, or reduction module satisfying the stated interface and memory restrictions.

\noindent\textbf{Role in the proof.}
Use in this chapter. It fixes the proof context, imported mathematical language, or quantified module parameters for the following statements.

\paragraph{Context 64 (line 630)}
\noindent\textbf{EasyCrypt name.}
\begin{lstlisting}[language=EasyCrypt,basicstyle=\ttfamily\scriptsize]
Bmlwe
\end{lstlisting}
\noindent\textbf{Checked EasyCrypt statement.}
\begin{lstlisting}[language=EasyCrypt,basicstyle=\ttfamily\scriptsize]
declare module Bmlwe <: MLWE_Reduction.
\end{lstlisting}
\noindent\textbf{Formal mathematical reading.}
Mathematical context. This is universal quantification over an adversary, oracle, scheme, sampler, or reduction module satisfying the stated interface and memory restrictions.

\noindent\textbf{Role in the proof.}
Use in this chapter. It fixes the proof context, imported mathematical language, or quantified module parameters for the following statements.

\paragraph{Context 65 (line 631)}
\noindent\textbf{EasyCrypt name.}
\begin{lstlisting}[language=EasyCrypt,basicstyle=\ttfamily\scriptsize]
Bbimodal
\end{lstlisting}
\noindent\textbf{Checked EasyCrypt statement.}
\begin{lstlisting}[language=EasyCrypt,basicstyle=\ttfamily\scriptsize]
declare module Bbimodal <: BimodalSelfTargetMSIS_Reduction.
\end{lstlisting}
\noindent\textbf{Formal mathematical reading.}
Mathematical context. This is universal quantification over an adversary, oracle, scheme, sampler, or reduction module satisfying the stated interface and memory restrictions.

\noindent\textbf{Role in the proof.}
Use in this chapter. It fixes the proof context, imported mathematical language, or quantified module parameters for the following statements.

\subsubsection*{Definitions, Carrier Sets, Operations, Games, and Procedures}
\begin{formaldefinition}[EasyCrypt line 21]
\noindent\textbf{EasyCrypt name.}
\begin{lstlisting}[language=EasyCrypt,basicstyle=\ttfamily\scriptsize]
correctness_obligation
\end{lstlisting}
\noindent\textbf{Checked EasyCrypt statement.}
\begin{lstlisting}[language=EasyCrypt,basicstyle=\ttfamily\scriptsize]
op correctness_obligation =
  forall (sd : seed) (coins : random_coins)
         (m : message) (ctx : context),
    verify_internal haetae_mode (keygen_internal haetae_mode sd).`1 m ctx
      (sign_internal haetae_mode
        (keygen_internal haetae_mode sd).`2 m ctx coins).
\end{lstlisting}
\noindent\textbf{Formal mathematical reading.}
Mathematical definition. This is a total EasyCrypt operation.  The exact displayed statement gives the domain, codomain, and defining equation when present.

\noindent\textbf{Role in the proof.}
Use in this chapter. It is part of the exact formal vocabulary used by subsequent lemmas and games.

\end{formaldefinition}

\subsubsection*{Lemmas, Theorems, Relational Equivalences, and Their Proof Dependencies}
\begin{formaltheorem}[EasyCrypt line 28]
\noindent\textbf{EasyCrypt name.}
\begin{lstlisting}[language=EasyCrypt,basicstyle=\ttfamily\scriptsize]
correctness_obligation_holds
\end{lstlisting}
\noindent\textbf{Checked EasyCrypt statement.}
\begin{lstlisting}[language=EasyCrypt,basicstyle=\ttfamily\scriptsize]
lemma correctness_obligation_holds : correctness_obligation.
\end{lstlisting}
\noindent\textbf{Formal mathematical reading.}
Mathematical theorem. This lemma is a universally quantified proposition over the variables shown in the EasyCrypt statement.

\noindent\textbf{Role in the proof.}
Use in this chapter. It proves a structural correctness fact needed by later extraction or verification arguments.

\noindent\textbf{Proof-script interpretation.}
Proof-script reading. The machine proof decomposes the theorem into rewrites, program-logic obligations, relational equivalences, and arithmetic side conditions.
\emph{Main tactics visible in the proof script:} \ec{rewrite}, \ec{apply}.
\emph{Named identifiers syntactically cited in the proof block:}
\begin{lstlisting}[language=EasyCrypt,basicstyle=\ttfamily\scriptsize]
coins
correctness_obligation
ctx
keygen_internal
sd
verify_internal_sign_internal
\end{lstlisting}

\end{formaltheorem}

\begin{formaltheorem}[EasyCrypt line 40]
\noindent\textbf{EasyCrypt name.}
\begin{lstlisting}[language=EasyCrypt,basicstyle=\ttfamily\scriptsize]
correctness_perfect
\end{lstlisting}
\noindent\textbf{Checked EasyCrypt statement.}
\begin{lstlisting}[language=EasyCrypt,basicstyle=\ttfamily\scriptsize]
lemma correctness_perfect &m :
  Pr[SIG.Correctness(H, HAETAE, A).main() @ &m : res] = 0%r.
\end{lstlisting}
\noindent\textbf{Formal mathematical reading.}
Mathematical theorem. This theorem has the form \(\Pr[G:E] = B\) is a game-probability statement.  It is one of the exact hops, bad-event bounds, or final advantage inequalities used in the reduction.

\noindent\textbf{Role in the proof.}
Use in this chapter. It proves a structural correctness fact needed by later extraction or verification arguments.

\noindent\textbf{Proof-script interpretation.}
Proof-script reading. The machine proof decomposes the theorem into rewrites, program-logic obligations, relational equivalences, and arithmetic side conditions.
\emph{Main tactics visible in the proof script:} \ec{byphoare}, \ec{hoare}, \ec{proc}, \ec{inline}, \ec{wp}, \ec{call}, \ec{rnd}, \ec{conseq}, \ec{rewrite}, \ec{split}, \ec{apply}.
\emph{Named identifiers syntactically cited in the proof block:}
\begin{lstlisting}[language=EasyCrypt,basicstyle=\ttfamily\scriptsize]
coins
keygen_internal
result
result0
sd
valid_signature_sign_internal
verify_internal
verify_norm_ok_current
\end{lstlisting}

\end{formaltheorem}

\begin{formaltheorem}[EasyCrypt line 81]
\noindent\textbf{EasyCrypt name.}
\begin{lstlisting}[language=EasyCrypt,basicstyle=\ttfamily\scriptsize]
euf_cma_security_no_signing
\end{lstlisting}
\noindent\textbf{Checked EasyCrypt statement.}
\begin{lstlisting}[language=EasyCrypt,basicstyle=\ttfamily\scriptsize]
lemma euf_cma_security_no_signing &m :
  Pr[SIG.UF_NMA(H, HAETAE, N).main() @ &m : res] <=
    Pr[MLWE_Game(Bmlwe).main() @ &m : res] +
    Pr[BimodalSelfTargetMSIS_Game(Bbimodal).main() @ &m : res] =>
  Pr[BimodalSelfTargetMSIS_Game(Bbimodal).main() @ &m : res] <=
    Pr[ModuleSIS_Game(Bsis).main() @ &m : res] +
    bimodal_to_msis_reduction_loss_term =>
  Pr[MLWE_Game(Bmlwe).main() @ &m : res] <= mlwe_hardness_term =>
  Pr[ModuleSIS_Game(Bsis).main() @ &m : res] <= module_sis_hardness_term =>
  Pr[SIG.EUF_CMA(H, HAETAE, SIG.NMA_As_CMA(N)).main() @ &m : res] <=
    haetae_euf_bound.
\end{lstlisting}
\noindent\textbf{Formal mathematical reading.}
Mathematical theorem. This theorem has the form \(\Pr[G:E] \le B\) is a game-probability statement.  It is one of the exact hops, bad-event bounds, or final advantage inequalities used in the reduction.

\noindent\textbf{Role in the proof.}
Use in this chapter. It contributes directly to an advantage bound, bad-event bound, or real-arithmetic composition step.

\noindent\textbf{Proof-script interpretation.}
Proof-script reading. The machine proof decomposes the theorem into rewrites, program-logic obligations, relational equivalences, and arithmetic side conditions.
\emph{Main tactics visible in the proof script:} \ec{rewrite}, \ec{have}, \ec{apply}.
\emph{Named identifiers syntactically cited in the proof block:}
\begin{lstlisting}[language=EasyCrypt,basicstyle=\ttfamily\scriptsize]
Bbimodal
BimodalSelfTargetMSIS_Game
Bmlwe
Bsis
HAETAE
MLWE_Game
ModuleSIS_Game
Pr
SIG
UF_NMA
bimodal_hop
bimodal_to_msis_reduction_loss_term
h_nma_msis
haetae_euf_bound_rejection_groupedE
ler_add
ler_addl
ler_trans
lerr
main
mlwe_bound
mlwe_hardness_term
module_sis_hardness_term
msis_bound
nma_as_cma_exact
nma_hop
rejection_bound_parts_nonnegative
rejection_sampling_bound_obligation_holds
\end{lstlisting}

\end{formaltheorem}

\begin{formaltheorem}[EasyCrypt line 136]
\noindent\textbf{EasyCrypt name.}
\begin{lstlisting}[language=EasyCrypt,basicstyle=\ttfamily\scriptsize]
euf_cma_security
\end{lstlisting}
\noindent\textbf{Checked EasyCrypt statement.}
\begin{lstlisting}[language=EasyCrypt,basicstyle=\ttfamily\scriptsize]
lemma euf_cma_security &m :
  Pr[SIG.EUF_CMA(H, HAETAE, A).main() @ &m : res] <=
    Pr[SIG.UF_NMA(H, HAETAE, N).main() @ &m : res] +
    (fs_with_aborts_reprogramming_term +
      fs_with_aborts_min_entropy_term) =>
  Pr[SIG.UF_NMA(H, HAETAE, N).main() @ &m : res] <=
    Pr[MLWE_Game(Bmlwe).main() @ &m : res] +
    Pr[BimodalSelfTargetMSIS_Game(Bbimodal).main() @ &m : res] =>
  Pr[BimodalSelfTargetMSIS_Game(Bbimodal).main() @ &m : res] <=
    Pr[ModuleSIS_Game(Bsis).main() @ &m : res] +
    bimodal_to_msis_reduction_loss_term =>
  Pr[MLWE_Game(Bmlwe).main() @ &m : res] <= mlwe_hardness_term =>
  Pr[ModuleSIS_Game(Bsis).main() @ &m : res] <= module_sis_hardness_term =>
  Pr[SIG.EUF_CMA(H, HAETAE, A).main() @ &m : res] <= haetae_euf_bound.
\end{lstlisting}
\noindent\textbf{Formal mathematical reading.}
Mathematical theorem. This theorem has the form \(\Pr[G:E] \le B\) is a game-probability statement.  It is one of the exact hops, bad-event bounds, or final advantage inequalities used in the reduction.

\noindent\textbf{Role in the proof.}
Use in this chapter. It contributes directly to an advantage bound, bad-event bound, or real-arithmetic composition step.

\noindent\textbf{Proof-script interpretation.}
Proof-script reading. The machine proof decomposes the theorem into rewrites, program-logic obligations, relational equivalences, and arithmetic side conditions.
\emph{Main tactics visible in the proof script:} \ec{have}, \ec{apply}, \ec{rewrite}.
\emph{Named identifiers syntactically cited in the proof block:}
\begin{lstlisting}[language=EasyCrypt,basicstyle=\ttfamily\scriptsize]
Bbimodal
BimodalSelfTargetMSIS_Game
Bmlwe
Bsis
HAETAE
MLWE_Game
ModuleSIS_Game
Pr
SIG
UF_NMA
bimodal_hop
bimodal_to_msis_reduction_loss_term
fs_hop
fs_with_aborts_min_entropy_term
fs_with_aborts_reprogramming_term
h_nma_msis
haetae_euf_bound_groupedE
ler_add
ler_trans
lerr
main
mlwe_bound
msis_bound
nma_hop
\end{lstlisting}

\end{formaltheorem}

\begin{formaltheorem}[EasyCrypt line 188]
\noindent\textbf{EasyCrypt name.}
\begin{lstlisting}[language=EasyCrypt,basicstyle=\ttfamily\scriptsize]
euf_cma_security_with_hop_interfaces
\end{lstlisting}
\noindent\textbf{Checked EasyCrypt statement.}
\begin{lstlisting}[language=EasyCrypt,basicstyle=\ttfamily\scriptsize]
lemma euf_cma_security_with_hop_interfaces &m :
  fs_with_aborts_interfaces_ready haetae_mode =>
  Pr[SIG.EUF_CMA(H, HAETAE, A).main() @ &m : res] <=
    Pr[SIG.UF_NMA(H, HAETAE, N).main() @ &m : res] +
    (fs_with_aborts_reprogramming_term +
      fs_with_aborts_min_entropy_term) =>
  Pr[SIG.UF_NMA(H, HAETAE, N).main() @ &m : res] <=
    Pr[MLWE_Game(Bmlwe).main() @ &m : res] +
    Pr[BimodalSelfTargetMSIS_Game(Bbimodal).main() @ &m : res] =>
  Pr[BimodalSelfTargetMSIS_Game(Bbimodal).main() @ &m : res] <=
    Pr[ModuleSIS_Game(Bsis).main() @ &m : res] +
    bimodal_to_msis_reduction_loss_term =>
  Pr[MLWE_Game(Bmlwe).main() @ &m : res] <= mlwe_hardness_term =>
  Pr[ModuleSIS_Game(Bsis).main() @ &m : res] <= module_sis_hardness_term =>
  Pr[SIG.EUF_CMA(H, HAETAE, A).main() @ &m : res] <= haetae_euf_bound.
\end{lstlisting}
\noindent\textbf{Formal mathematical reading.}
Mathematical theorem. This theorem has the form \(\Pr[G:E] \le B\) is a game-probability statement.  It is one of the exact hops, bad-event bounds, or final advantage inequalities used in the reduction.

\noindent\textbf{Role in the proof.}
Use in this chapter. It contributes directly to an advantage bound, bad-event bound, or real-arithmetic composition step.

\noindent\textbf{Proof-script interpretation.}
Proof-script reading. The machine proof decomposes the theorem into rewrites, program-logic obligations, relational equivalences, and arithmetic side conditions.
\emph{Main tactics visible in the proof script:} \ec{apply}.
\emph{Named identifiers syntactically cited in the proof block:}
\begin{lstlisting}[language=EasyCrypt,basicstyle=\ttfamily\scriptsize]
bimodal_hop
euf_cma_security
fs_hop
mlwe_bound
msis_bound
nma_hop
\end{lstlisting}

\end{formaltheorem}

\begin{formaltheorem}[EasyCrypt line 223]
\noindent\textbf{EasyCrypt name.}
\begin{lstlisting}[language=EasyCrypt,basicstyle=\ttfamily\scriptsize]
euf_cma_security_from_simulated_hop
\end{lstlisting}
\noindent\textbf{Checked EasyCrypt statement.}
\begin{lstlisting}[language=EasyCrypt,basicstyle=\ttfamily\scriptsize]
lemma euf_cma_security_from_simulated_hop &m :
  fs_with_aborts_interfaces_ready haetae_mode =>
  Pr[EUF_CMA_SimulatedSign(H, HAETAE, A, RealSigningSimulator(HAETAE)).main()
       @ &m : res] <=
    Pr[SIG.UF_NMA(H, HAETAE, N).main() @ &m : res] +
    (fs_with_aborts_reprogramming_term +
      fs_with_aborts_min_entropy_term) =>
  Pr[SIG.UF_NMA(H, HAETAE, N).main() @ &m : res] <=
    Pr[MLWE_Game(Bmlwe).main() @ &m : res] +
    Pr[BimodalSelfTargetMSIS_Game(Bbimodal).main() @ &m : res] =>
  Pr[BimodalSelfTargetMSIS_Game(Bbimodal).main() @ &m : res] <=
    Pr[ModuleSIS_Game(Bsis).main() @ &m : res] +
    bimodal_to_msis_reduction_loss_term =>
  Pr[MLWE_Game(Bmlwe).main() @ &m : res] <= mlwe_hardness_term =>
  Pr[ModuleSIS_Game(Bsis).main() @ &m : res] <= module_sis_hardness_term =>
  Pr[SIG.EUF_CMA(H, HAETAE, A).main() @ &m : res] <= haetae_euf_bound.
\end{lstlisting}
\noindent\textbf{Formal mathematical reading.}
Mathematical theorem. This theorem has the form \(\Pr[G:E] \le B\) is a game-probability statement.  It is one of the exact hops, bad-event bounds, or final advantage inequalities used in the reduction.

\noindent\textbf{Role in the proof.}
Use in this chapter. It contributes directly to an advantage bound, bad-event bound, or real-arithmetic composition step.

\noindent\textbf{Proof-script interpretation.}
Proof-script reading. The machine proof decomposes the theorem into rewrites, program-logic obligations, relational equivalences, and arithmetic side conditions.
\emph{Main tactics visible in the proof script:} \ec{rewrite}, \ec{have}, \ec{apply}.
\emph{Named identifiers syntactically cited in the proof block:}
\begin{lstlisting}[language=EasyCrypt,basicstyle=\ttfamily\scriptsize]
Bbimodal
BimodalSelfTargetMSIS_Game
Bmlwe
Bsis
HAETAE
MLWE_Game
ModuleSIS_Game
Pr
SIG
UF_NMA
bimodal_hop
bimodal_to_msis_reduction_loss_term
fs_with_aborts_min_entropy_term
fs_with_aborts_reprogramming_term
h_nma_msis
haetae_euf_bound_groupedE
iface
ler_add
ler_trans
lerr
main
mlwe_bound
msis_bound
nma_hop
real_signing_simulator_exact
sim_hop
\end{lstlisting}

\end{formaltheorem}

\begin{formaltheorem}[EasyCrypt line 278]
\noindent\textbf{EasyCrypt name.}
\begin{lstlisting}[language=EasyCrypt,basicstyle=\ttfamily\scriptsize]
euf_cma_security_from_explicit_boundaries
\end{lstlisting}
\noindent\textbf{Checked EasyCrypt statement.}
\begin{lstlisting}[language=EasyCrypt,basicstyle=\ttfamily\scriptsize]
lemma euf_cma_security_from_explicit_boundaries &m :
  Pr[EUF_CMA_SimulatedSign(H, HAETAE, A, RealSigningSimulator(HAETAE)).main()
       @ &m : res] <=
    Pr[SIG.UF_NMA(H, HAETAE, N).main() @ &m : res] +
    (fs_with_aborts_reprogramming_term +
      fs_with_aborts_min_entropy_term) =>
  Pr[SIG.UF_NMA(H, HAETAE, N).main() @ &m : res] <=
    Pr[MLWE_Game(Bmlwe).main() @ &m : res] +
    Pr[BimodalSelfTargetMSIS_Game(Bbimodal).main() @ &m : res] =>
  Pr[BimodalSelfTargetMSIS_Game(Bbimodal).main() @ &m : res] <=
    Pr[ModuleSIS_Game(Bsis).main() @ &m : res] +
    bimodal_to_msis_reduction_loss_term =>
  Pr[MLWE_Game(Bmlwe).main() @ &m : res] <= mlwe_hardness_term =>
  Pr[ModuleSIS_Game(Bsis).main() @ &m : res] <= module_sis_hardness_term =>
  Pr[SIG.EUF_CMA(H, HAETAE, A).main() @ &m : res] <= haetae_euf_bound.
\end{lstlisting}
\noindent\textbf{Formal mathematical reading.}
Mathematical theorem. This theorem has the form \(\Pr[G:E] \le B\) is a game-probability statement.  It is one of the exact hops, bad-event bounds, or final advantage inequalities used in the reduction.

\noindent\textbf{Role in the proof.}
Use in this chapter. It contributes directly to an advantage bound, bad-event bound, or real-arithmetic composition step.

\noindent\textbf{Proof-script interpretation.}
Proof-script reading. The machine proof decomposes the theorem into rewrites, program-logic obligations, relational equivalences, and arithmetic side conditions.
\emph{Main tactics visible in the proof script:} \ec{apply}.
\emph{Named identifiers syntactically cited in the proof block:}
\begin{lstlisting}[language=EasyCrypt,basicstyle=\ttfamily\scriptsize]
bimodal_hop
euf_cma_security_from_simulated_hop
fs_with_aborts_interfaces_ready_holds
mlwe_bound
msis_bound
nma_hop
sim_hop
\end{lstlisting}

\end{formaltheorem}

\begin{formaltheorem}[EasyCrypt line 319]
\noindent\textbf{EasyCrypt name.}
\begin{lstlisting}[language=EasyCrypt,basicstyle=\ttfamily\scriptsize]
euf_cma_security_from_rom_internal_transcript_hop
\end{lstlisting}
\noindent\textbf{Checked EasyCrypt statement.}
\begin{lstlisting}[language=EasyCrypt,basicstyle=\ttfamily\scriptsize]
lemma euf_cma_security_from_rom_internal_transcript_hop &m :
  Pr[EUF_CMA_ROMInternalTranscriptSign(H, A).main() @ &m : res] <=
    Pr[SIG.UF_NMA(H, HAETAE, N).main() @ &m : res] +
    (fs_with_aborts_reprogramming_term +
      fs_with_aborts_min_entropy_term) =>
  Pr[SIG.UF_NMA(H, HAETAE, N).main() @ &m : res] <=
    Pr[MLWE_Game(Bmlwe).main() @ &m : res] +
    Pr[BimodalSelfTargetMSIS_Game(Bbimodal).main() @ &m : res] =>
  Pr[BimodalSelfTargetMSIS_Game(Bbimodal).main() @ &m : res] <=
    Pr[ModuleSIS_Game(Bsis).main() @ &m : res] +
    bimodal_to_msis_reduction_loss_term =>
  Pr[MLWE_Game(Bmlwe).main() @ &m : res] <= mlwe_hardness_term =>
  Pr[ModuleSIS_Game(Bsis).main() @ &m : res] <= module_sis_hardness_term =>
  Pr[SIG.EUF_CMA(H, HAETAE, A).main() @ &m : res] <= haetae_euf_bound.
\end{lstlisting}
\noindent\textbf{Formal mathematical reading.}
Mathematical theorem. This theorem has the form \(\Pr[G:E] \le B\) is a game-probability statement.  It is one of the exact hops, bad-event bounds, or final advantage inequalities used in the reduction.

\noindent\textbf{Role in the proof.}
Use in this chapter. It contributes directly to an advantage bound, bad-event bound, or real-arithmetic composition step.

\noindent\textbf{Proof-script interpretation.}
Proof-script reading. The machine proof decomposes the theorem into rewrites, program-logic obligations, relational equivalences, and arithmetic side conditions.
\emph{Main tactics visible in the proof script:} \ec{apply}, \ec{rewrite}.
\emph{Named identifiers syntactically cited in the proof block:}
\begin{lstlisting}[language=EasyCrypt,basicstyle=\ttfamily\scriptsize]
Bbimodal
Bmlwe
Bsis
bimodal_hop
euf_cma_security
haetae_rom_internal_transcript_erasure_exact
mlwe_bound
msis_bound
nma_hop
rom_hop
\end{lstlisting}

\end{formaltheorem}

\begin{formaltheorem}[EasyCrypt line 344]
\noindent\textbf{EasyCrypt name.}
\begin{lstlisting}[language=EasyCrypt,basicstyle=\ttfamily\scriptsize]
rom_internal_transcript_hop_from_clear_site_reduction
\end{lstlisting}
\noindent\textbf{Checked EasyCrypt statement.}
\begin{lstlisting}[language=EasyCrypt,basicstyle=\ttfamily\scriptsize]
lemma rom_internal_transcript_hop_from_clear_site_reduction &m :
  Pr[EUF_CMA_ROMInternalTranscriptSign(H, A).main() @ &m :
      res /\
      transcript_log_signature_programming_sites_clear haetae_mode
        EUF_CMA_ROMInternalTranscriptSign.transcripts] <=
    Pr[SIG.UF_NMA(H, HAETAE, N).main() @ &m : res] =>
  Pr[EUF_CMA_ROMInternalTranscriptSign(H, A).main() @ &m : res] <=
    Pr[SIG.UF_NMA(H, HAETAE, N).main() @ &m : res] +
    (fs_with_aborts_reprogramming_term +
      fs_with_aborts_min_entropy_term).
\end{lstlisting}
\noindent\textbf{Formal mathematical reading.}
Mathematical theorem. This theorem has the form \(\Pr[G:E] \le B\) is a game-probability statement.  It is one of the exact hops, bad-event bounds, or final advantage inequalities used in the reduction.

\noindent\textbf{Role in the proof.}
Use in this chapter. It contributes directly to an advantage bound, bad-event bound, or real-arithmetic composition step.

\noindent\textbf{Proof-script interpretation.}
Proof-script reading. The machine proof decomposes the theorem into rewrites, program-logic obligations, relational equivalences, and arithmetic side conditions.
\emph{Main tactics visible in the proof script:} \ec{have}, \ec{rewrite}, \ec{apply}.
\emph{Named identifiers syntactically cited in the proof block:}
\begin{lstlisting}[language=EasyCrypt,basicstyle=\ttfamily\scriptsize]
EUF_CMA_ROMInternalTranscriptSign
Pr
clear_reduction
haetae_mode
ler_add
main
mu_split
rom_internal_transcript_bad_forgery_bound
split_success
transcript_log_signature_programming_sites_clear
transcripts
\end{lstlisting}

\end{formaltheorem}

\begin{formaltheorem}[EasyCrypt line 375]
\noindent\textbf{EasyCrypt name.}
\begin{lstlisting}[language=EasyCrypt,basicstyle=\ttfamily\scriptsize]
rom_internal_transcript_hop_from_clear_site_counted_reduction
\end{lstlisting}
\noindent\textbf{Checked EasyCrypt statement.}
\begin{lstlisting}[language=EasyCrypt,basicstyle=\ttfamily\scriptsize]
lemma rom_internal_transcript_hop_from_clear_site_counted_reduction &m :
  Pr[EUF_CMA_ROMInternalTranscriptSign(H, A).main() @ &m :
      res /\
      transcript_log_signature_programming_sites_clear haetae_mode
        EUF_CMA_ROMInternalTranscriptSign.transcripts] <=
    Pr[SIG.UF_NMA(H, HAETAE, N).main() @ &m : res] =>
  Pr[EUF_CMA_ROMInternalTranscriptSign(H, A).main() @ &m : res] <=
    Pr[SIG.UF_NMA(H, HAETAE, N).main() @ &m : res] +
    counted_rom_programming_loss_term.
\end{lstlisting}
\noindent\textbf{Formal mathematical reading.}
Mathematical theorem. This theorem has the form \(\Pr[G:E] \le B\) is a game-probability statement.  It is one of the exact hops, bad-event bounds, or final advantage inequalities used in the reduction.

\noindent\textbf{Role in the proof.}
Use in this chapter. It contributes directly to an advantage bound, bad-event bound, or real-arithmetic composition step.

\noindent\textbf{Proof-script interpretation.}
Proof-script reading. The machine proof decomposes the theorem into rewrites, program-logic obligations, relational equivalences, and arithmetic side conditions.
\emph{Main tactics visible in the proof script:} \ec{have}, \ec{rewrite}, \ec{apply}.
\emph{Named identifiers syntactically cited in the proof block:}
\begin{lstlisting}[language=EasyCrypt,basicstyle=\ttfamily\scriptsize]
EUF_CMA_ROMInternalTranscriptSign
Pr
clear_reduction
haetae_mode
ler_add
main
mu_split
rom_internal_transcript_bad_forgery_counted_bound
split_success
transcript_log_signature_programming_sites_clear
transcripts
\end{lstlisting}

\end{formaltheorem}

\begin{formaltheorem}[EasyCrypt line 418]
\noindent\textbf{EasyCrypt name.}
\begin{lstlisting}[language=EasyCrypt,basicstyle=\ttfamily\scriptsize]
bimodal_to_module_sis_reduction_bound
\end{lstlisting}
\noindent\textbf{Checked EasyCrypt statement.}
\begin{lstlisting}[language=EasyCrypt,basicstyle=\ttfamily\scriptsize]
lemma bimodal_to_module_sis_reduction_bound &m :
  Pr[BimodalSelfTargetMSIS_Game(Bbimodal).main() @ &m : res] <=
  Pr[ModuleSIS_Game(BimodalMSIS_As_ModuleSIS(Bbimodal)).main()
       @ &m : res] +
    bimodal_to_msis_reduction_loss_term.
\end{lstlisting}
\noindent\textbf{Formal mathematical reading.}
Mathematical theorem. This theorem has the form \(\Pr[G:E] \le B\) is a game-probability statement.  It is one of the exact hops, bad-event bounds, or final advantage inequalities used in the reduction.

\noindent\textbf{Role in the proof.}
Use in this chapter. It contributes directly to an advantage bound, bad-event bound, or real-arithmetic composition step.

\noindent\textbf{Proof-script interpretation.}
Proof-script reading. The machine proof decomposes the theorem into rewrites, program-logic obligations, relational equivalences, and arithmetic side conditions.
\emph{Main tactics visible in the proof script:} \ec{rewrite}, \ec{have}.
\emph{Named identifiers syntactically cited in the proof block:}
\begin{lstlisting}[language=EasyCrypt,basicstyle=\ttfamily\scriptsize]
Bbimodal
BimodalMSIS_As_ModuleSIS
ModuleSIS_Game
Pr
bimodal_msis_as_module_sis_exact
bimodal_to_msis_reduction_loss_termE
lerr
main
ring
\end{lstlisting}

\end{formaltheorem}

\begin{formaltheorem}[EasyCrypt line 434]
\noindent\textbf{EasyCrypt name.}
\begin{lstlisting}[language=EasyCrypt,basicstyle=\ttfamily\scriptsize]
euf_cma_security_from_mechanized_boundaries
\end{lstlisting}
\noindent\textbf{Checked EasyCrypt statement.}
\begin{lstlisting}[language=EasyCrypt,basicstyle=\ttfamily\scriptsize]
lemma euf_cma_security_from_mechanized_boundaries &m :
  Pr[EUF_CMA_SimulatedSign(H, HAETAE, A, RealSigningSimulator(HAETAE)).main()
       @ &m : res] <=
    Pr[SIG.UF_NMA(H, HAETAE, N).main() @ &m : res] +
    (fs_with_aborts_reprogramming_term +
      fs_with_aborts_min_entropy_term) =>
  Pr[SIG.UF_NMA(H, HAETAE, N).main() @ &m : res] <=
    Pr[MLWE_Game(Bmlwe).main() @ &m : res] +
    Pr[BimodalSelfTargetMSIS_Game(Bbimodal).main() @ &m : res] =>
  Pr[MLWE_Game(Bmlwe).main() @ &m : res] <= mlwe_hardness_term =>
  Pr[ModuleSIS_Game(BimodalMSIS_As_ModuleSIS(Bbimodal)).main()
       @ &m : res] <= module_sis_hardness_term =>
  Pr[SIG.EUF_CMA(H, HAETAE, A).main() @ &m : res] <= haetae_euf_bound.
\end{lstlisting}
\noindent\textbf{Formal mathematical reading.}
Mathematical theorem. This theorem has the form \(\Pr[G:E] \le B\) is a game-probability statement.  It is one of the exact hops, bad-event bounds, or final advantage inequalities used in the reduction.

\noindent\textbf{Role in the proof.}
Use in this chapter. It contributes directly to an advantage bound, bad-event bound, or real-arithmetic composition step.

\noindent\textbf{Proof-script interpretation.}
Proof-script reading. The machine proof decomposes the theorem into rewrites, program-logic obligations, relational equivalences, and arithmetic side conditions.
\emph{Main tactics visible in the proof script:} \ec{apply}.
\emph{Named identifiers syntactically cited in the proof block:}
\begin{lstlisting}[language=EasyCrypt,basicstyle=\ttfamily\scriptsize]
Bbimodal
BimodalMSIS_As_ModuleSIS
Bmlwe
bimodal_to_module_sis_reduction_bound
euf_cma_security_from_simulated_hop
fs_with_aborts_interfaces_ready_holds
mlwe_bound
msis_bound
nma_hop
sim_hop
\end{lstlisting}

\end{formaltheorem}

\begin{formaltheorem}[EasyCrypt line 473]
\noindent\textbf{EasyCrypt name.}
\begin{lstlisting}[language=EasyCrypt,basicstyle=\ttfamily\scriptsize]
euf_cma_security_from_rom_transcript_and_mechanized_bimodal
\end{lstlisting}
\noindent\textbf{Checked EasyCrypt statement.}
\begin{lstlisting}[language=EasyCrypt,basicstyle=\ttfamily\scriptsize]
lemma euf_cma_security_from_rom_transcript_and_mechanized_bimodal &m :
  Pr[EUF_CMA_ROMInternalTranscriptSign(H, A).main() @ &m : res] <=
    Pr[SIG.UF_NMA(H, HAETAE, N).main() @ &m : res] +
    (fs_with_aborts_reprogramming_term +
      fs_with_aborts_min_entropy_term) =>
  Pr[SIG.UF_NMA(H, HAETAE, N).main() @ &m : res] <=
    Pr[MLWE_Game(Bmlwe).main() @ &m : res] +
    Pr[BimodalSelfTargetMSIS_Game(Bbimodal).main() @ &m : res] =>
  Pr[MLWE_Game(Bmlwe).main() @ &m : res] <= mlwe_hardness_term =>
  Pr[ModuleSIS_Game(BimodalMSIS_As_ModuleSIS(Bbimodal)).main()
       @ &m : res] <= module_sis_hardness_term =>
  Pr[SIG.EUF_CMA(H, HAETAE, A).main() @ &m : res] <= haetae_euf_bound.
\end{lstlisting}
\noindent\textbf{Formal mathematical reading.}
Mathematical theorem. This theorem has the form \(\Pr[G:E] \le B\) is a game-probability statement.  It is one of the exact hops, bad-event bounds, or final advantage inequalities used in the reduction.

\noindent\textbf{Role in the proof.}
Use in this chapter. It contributes directly to an advantage bound, bad-event bound, or real-arithmetic composition step.

\noindent\textbf{Proof-script interpretation.}
Proof-script reading. The machine proof decomposes the theorem into rewrites, program-logic obligations, relational equivalences, and arithmetic side conditions.
\emph{Main tactics visible in the proof script:} \ec{apply}.
\emph{Named identifiers syntactically cited in the proof block:}
\begin{lstlisting}[language=EasyCrypt,basicstyle=\ttfamily\scriptsize]
Bbimodal
BimodalMSIS_As_ModuleSIS
Bmlwe
bimodal_to_module_sis_reduction_bound
euf_cma_security_from_rom_internal_transcript_hop
mlwe_bound
msis_bound
nma_hop
rom_hop
\end{lstlisting}

\end{formaltheorem}

\begin{formaltheorem}[EasyCrypt line 496]
\noindent\textbf{EasyCrypt name.}
\begin{lstlisting}[language=EasyCrypt,basicstyle=\ttfamily\scriptsize]
euf_cma_security_from_rom_transcript_and_counted_bimodal
\end{lstlisting}
\noindent\textbf{Checked EasyCrypt statement.}
\begin{lstlisting}[language=EasyCrypt,basicstyle=\ttfamily\scriptsize]
lemma euf_cma_security_from_rom_transcript_and_counted_bimodal &m :
  Pr[EUF_CMA_ROMInternalTranscriptSign(H, A).main() @ &m : res] <=
    Pr[SIG.UF_NMA(H, HAETAE, N).main() @ &m : res] +
    counted_rom_programming_loss_term =>
  Pr[SIG.UF_NMA(H, HAETAE, N).main() @ &m : res] <=
    Pr[MLWE_Game(Bmlwe).main() @ &m : res] +
    Pr[BimodalSelfTargetMSIS_Game(Bbimodal).main() @ &m : res] =>
  Pr[MLWE_Game(Bmlwe).main() @ &m : res] <= mlwe_hardness_term =>
  Pr[ModuleSIS_Game(BimodalMSIS_As_ModuleSIS(Bbimodal)).main()
       @ &m : res] <= module_sis_hardness_term =>
  Pr[SIG.EUF_CMA(H, HAETAE, A).main() @ &m : res] <=
    haetae_euf_counted_rom_bound.
\end{lstlisting}
\noindent\textbf{Formal mathematical reading.}
Mathematical theorem. This theorem has the form \(\Pr[G:E] \le B\) is a game-probability statement.  It is one of the exact hops, bad-event bounds, or final advantage inequalities used in the reduction.

\noindent\textbf{Role in the proof.}
Use in this chapter. It contributes directly to an advantage bound, bad-event bound, or real-arithmetic composition step.

\noindent\textbf{Proof-script interpretation.}
Proof-script reading. The machine proof decomposes the theorem into rewrites, program-logic obligations, relational equivalences, and arithmetic side conditions.
\emph{Main tactics visible in the proof script:} \ec{rewrite}, \ec{have}, \ec{apply}.
\emph{Named identifiers syntactically cited in the proof block:}
\begin{lstlisting}[language=EasyCrypt,basicstyle=\ttfamily\scriptsize]
Bbimodal
BimodalMSIS_As_ModuleSIS
BimodalSelfTargetMSIS_Game
Bmlwe
HAETAE
MLWE_Game
ModuleSIS_Game
Pr
SIG
UF_NMA
bimodal_to_module_sis_reduction_bound
bimodal_to_msis_reduction_loss_term
counted_rom_programming_loss_term
h_nma_msis
haetae_euf_counted_rom_bound_groupedE
haetae_rom_internal_transcript_erasure_exact
ler_add
ler_trans
lerr
main
mlwe_bound
msis_bound
nma_hop
rom_hop
\end{lstlisting}

\end{formaltheorem}

\begin{formaltheorem}[EasyCrypt line 547]
\noindent\textbf{EasyCrypt name.}
\begin{lstlisting}[language=EasyCrypt,basicstyle=\ttfamily\scriptsize]
euf_cma_security_from_rom_clear_site_and_mechanized_bimodal
\end{lstlisting}
\noindent\textbf{Checked EasyCrypt statement.}
\begin{lstlisting}[language=EasyCrypt,basicstyle=\ttfamily\scriptsize]
lemma euf_cma_security_from_rom_clear_site_and_mechanized_bimodal &m :
  Pr[EUF_CMA_ROMInternalTranscriptSign(H, A).main() @ &m :
      res /\
      transcript_log_signature_programming_sites_clear haetae_mode
        EUF_CMA_ROMInternalTranscriptSign.transcripts] <=
    Pr[SIG.UF_NMA(H, HAETAE, N).main() @ &m : res] =>
  Pr[SIG.UF_NMA(H, HAETAE, N).main() @ &m : res] <=
    Pr[MLWE_Game(Bmlwe).main() @ &m : res] +
    Pr[BimodalSelfTargetMSIS_Game(Bbimodal).main() @ &m : res] =>
  Pr[MLWE_Game(Bmlwe).main() @ &m : res] <= mlwe_hardness_term =>
  Pr[ModuleSIS_Game(BimodalMSIS_As_ModuleSIS(Bbimodal)).main()
       @ &m : res] <= module_sis_hardness_term =>
  Pr[SIG.EUF_CMA(H, HAETAE, A).main() @ &m : res] <= haetae_euf_bound.
\end{lstlisting}
\noindent\textbf{Formal mathematical reading.}
Mathematical theorem. This theorem has the form \(\Pr[G:E] \le B\) is a game-probability statement.  It is one of the exact hops, bad-event bounds, or final advantage inequalities used in the reduction.

\noindent\textbf{Role in the proof.}
Use in this chapter. It contributes directly to an advantage bound, bad-event bound, or real-arithmetic composition step.

\noindent\textbf{Proof-script interpretation.}
Proof-script reading. The machine proof decomposes the theorem into rewrites, program-logic obligations, relational equivalences, and arithmetic side conditions.
\emph{Main tactics visible in the proof script:} \ec{apply}.
\emph{Named identifiers syntactically cited in the proof block:}
\begin{lstlisting}[language=EasyCrypt,basicstyle=\ttfamily\scriptsize]
clear_reduction
euf_cma_security_from_rom_transcript_and_mechanized_bimodal
mlwe_bound
msis_bound
nma_hop
rom_internal_transcript_hop_from_clear_site_reduction
\end{lstlisting}

\end{formaltheorem}

\begin{formaltheorem}[EasyCrypt line 571]
\noindent\textbf{EasyCrypt name.}
\begin{lstlisting}[language=EasyCrypt,basicstyle=\ttfamily\scriptsize]
euf_cma_security_from_rom_clear_site_and_counted_bimodal
\end{lstlisting}
\noindent\textbf{Checked EasyCrypt statement.}
\begin{lstlisting}[language=EasyCrypt,basicstyle=\ttfamily\scriptsize]
lemma euf_cma_security_from_rom_clear_site_and_counted_bimodal &m :
  Pr[EUF_CMA_ROMInternalTranscriptSign(H, A).main() @ &m :
      res /\
      transcript_log_signature_programming_sites_clear haetae_mode
        EUF_CMA_ROMInternalTranscriptSign.transcripts] <=
    Pr[SIG.UF_NMA(H, HAETAE, N).main() @ &m : res] =>
  Pr[SIG.UF_NMA(H, HAETAE, N).main() @ &m : res] <=
    Pr[MLWE_Game(Bmlwe).main() @ &m : res] +
    Pr[BimodalSelfTargetMSIS_Game(Bbimodal).main() @ &m : res] =>
  Pr[MLWE_Game(Bmlwe).main() @ &m : res] <= mlwe_hardness_term =>
  Pr[ModuleSIS_Game(BimodalMSIS_As_ModuleSIS(Bbimodal)).main()
       @ &m : res] <= module_sis_hardness_term =>
  Pr[SIG.EUF_CMA(H, HAETAE, A).main() @ &m : res] <=
    haetae_euf_counted_rom_bound.
\end{lstlisting}
\noindent\textbf{Formal mathematical reading.}
Mathematical theorem. This theorem has the form \(\Pr[G:E] \le B\) is a game-probability statement.  It is one of the exact hops, bad-event bounds, or final advantage inequalities used in the reduction.

\noindent\textbf{Role in the proof.}
Use in this chapter. It contributes directly to an advantage bound, bad-event bound, or real-arithmetic composition step.

\noindent\textbf{Proof-script interpretation.}
Proof-script reading. The machine proof decomposes the theorem into rewrites, program-logic obligations, relational equivalences, and arithmetic side conditions.
\emph{Main tactics visible in the proof script:} \ec{apply}.
\emph{Named identifiers syntactically cited in the proof block:}
\begin{lstlisting}[language=EasyCrypt,basicstyle=\ttfamily\scriptsize]
clear_reduction
euf_cma_security_from_rom_transcript_and_counted_bimodal
mlwe_bound
msis_bound
nma_hop
rom_internal_transcript_hop_from_clear_site_counted_reduction
\end{lstlisting}

\end{formaltheorem}

\begin{formaltheorem}[EasyCrypt line 633]
\noindent\textbf{EasyCrypt name.}
\begin{lstlisting}[language=EasyCrypt,basicstyle=\ttfamily\scriptsize]
euf_cma_security_from_structural_nma_and_counted_bimodal
\end{lstlisting}
\noindent\textbf{Checked EasyCrypt statement.}
\begin{lstlisting}[language=EasyCrypt,basicstyle=\ttfamily\scriptsize]
lemma euf_cma_security_from_structural_nma_and_counted_bimodal &m :
  Pr[SIG.UF_NMA(H, HAETAE, ROMInternalTranscriptAsNMA(A)).main()
       @ &m : res] <=
    Pr[MLWE_Game(Bmlwe).main() @ &m : res] +
    Pr[BimodalSelfTargetMSIS_Game(Bbimodal).main() @ &m : res] =>
  Pr[MLWE_Game(Bmlwe).main() @ &m : res] <= mlwe_hardness_term =>
  Pr[ModuleSIS_Game(BimodalMSIS_As_ModuleSIS(Bbimodal)).main()
       @ &m : res] <= module_sis_hardness_term =>
  Pr[SIG.EUF_CMA(H, HAETAE, A).main() @ &m : res] <=
    haetae_euf_counted_rom_bound.
\end{lstlisting}
\noindent\textbf{Formal mathematical reading.}
Mathematical theorem. This theorem has the form \(\Pr[G:E] \le B\) is a game-probability statement.  It is one of the exact hops, bad-event bounds, or final advantage inequalities used in the reduction.

\noindent\textbf{Role in the proof.}
Use in this chapter. It contributes directly to an advantage bound, bad-event bound, or real-arithmetic composition step.

\noindent\textbf{Proof-script interpretation.}
Proof-script reading. The machine proof decomposes the theorem into rewrites, program-logic obligations, relational equivalences, and arithmetic side conditions.
\emph{Main tactics visible in the proof script:} \ec{apply}.
\emph{Named identifiers syntactically cited in the proof block:}
\begin{lstlisting}[language=EasyCrypt,basicstyle=\ttfamily\scriptsize]
Bbimodal
Bmlwe
ROMInternalTranscriptAsNMA
euf_cma_security_from_rom_clear_site_and_counted_bimodal
mlwe_bound
msis_bound
nma_hop
rom_internal_transcript_clear_site_structural_nma_bound
\end{lstlisting}

\end{formaltheorem}

\begin{formaltheorem}[EasyCrypt line 653]
\noindent\textbf{EasyCrypt name.}
\begin{lstlisting}[language=EasyCrypt,basicstyle=\ttfamily\scriptsize]
euf_cma_security_from_public_sim_nma_and_counted_bimodal
\end{lstlisting}
\noindent\textbf{Checked EasyCrypt statement.}
\begin{lstlisting}[language=EasyCrypt,basicstyle=\ttfamily\scriptsize]
lemma euf_cma_security_from_public_sim_nma_and_counted_bimodal &m :
  Pr[SIG.UF_NMA(H, HAETAE, ROMInternalTranscriptPublicSimAsNMA(A)).main()
       @ &m : res] <=
    Pr[MLWE_Game(Bmlwe).main() @ &m : res] +
    Pr[BimodalSelfTargetMSIS_Game(Bbimodal).main() @ &m : res] =>
  Pr[MLWE_Game(Bmlwe).main() @ &m : res] <= mlwe_hardness_term =>
  Pr[ModuleSIS_Game(BimodalMSIS_As_ModuleSIS(Bbimodal)).main()
       @ &m : res] <= module_sis_hardness_term =>
  Pr[SIG.EUF_CMA(H, HAETAE, A).main() @ &m : res] <=
    haetae_euf_counted_rom_bound.
\end{lstlisting}
\noindent\textbf{Formal mathematical reading.}
Mathematical theorem. This theorem has the form \(\Pr[G:E] \le B\) is a game-probability statement.  It is one of the exact hops, bad-event bounds, or final advantage inequalities used in the reduction.

\noindent\textbf{Role in the proof.}
Use in this chapter. It contributes directly to an advantage bound, bad-event bound, or real-arithmetic composition step.

\noindent\textbf{Proof-script interpretation.}
Proof-script reading. The machine proof decomposes the theorem into rewrites, program-logic obligations, relational equivalences, and arithmetic side conditions.
\emph{Main tactics visible in the proof script:} \ec{apply}, \ec{rewrite}.
\emph{Named identifiers syntactically cited in the proof block:}
\begin{lstlisting}[language=EasyCrypt,basicstyle=\ttfamily\scriptsize]
euf_cma_security_from_structural_nma_and_counted_bimodal
mlwe_bound
msis_bound
nma_hop
rom_internal_nma_public_sim_exact
\end{lstlisting}

\end{formaltheorem}

\begin{formaltheorem}[EasyCrypt line 673]
\noindent\textbf{EasyCrypt name.}
\begin{lstlisting}[language=EasyCrypt,basicstyle=\ttfamily\scriptsize]
euf_cma_security_from_paper_sim_nma_and_counted_bimodal
\end{lstlisting}
\noindent\textbf{Checked EasyCrypt statement.}
\begin{lstlisting}[language=EasyCrypt,basicstyle=\ttfamily\scriptsize]
lemma euf_cma_security_from_paper_sim_nma_and_counted_bimodal &m :
  Pr[SIG.UF_NMA(H, HAETAE, ROMInternalTranscriptPublicSimAsNMA(A)).main()
       @ &m : res] <=
  Pr[SIG.UF_NMA(H, HAETAE,
       ROMInternalTranscriptPaperSimAsNMA(A, Samp)).main() @ &m : res] +
    rejection_sampling_loss_term =>
  Pr[SIG.UF_NMA(H, HAETAE,
       ROMInternalTranscriptPaperSimAsNMA(A, Samp)).main() @ &m : res] +
    rejection_sampling_loss_term <=
    Pr[MLWE_Game(Bmlwe).main() @ &m : res] +
    Pr[BimodalSelfTargetMSIS_Game(Bbimodal).main() @ &m : res] =>
  Pr[MLWE_Game(Bmlwe).main() @ &m : res] <= mlwe_hardness_term =>
  Pr[ModuleSIS_Game(BimodalMSIS_As_ModuleSIS(Bbimodal)).main()
       @ &m : res] <= module_sis_hardness_term =>
  Pr[SIG.EUF_CMA(H, HAETAE, A).main() @ &m : res] <=
    haetae_euf_counted_rom_bound.
\end{lstlisting}
\noindent\textbf{Formal mathematical reading.}
Mathematical theorem. This theorem has the form \(\Pr[G:E] \le B\) is a game-probability statement.  It is one of the exact hops, bad-event bounds, or final advantage inequalities used in the reduction.

\noindent\textbf{Role in the proof.}
Use in this chapter. It contributes directly to an advantage bound, bad-event bound, or real-arithmetic composition step.

\noindent\textbf{Proof-script interpretation.}
Proof-script reading. The machine proof decomposes the theorem into rewrites, program-logic obligations, relational equivalences, and arithmetic side conditions.
\emph{Main tactics visible in the proof script:} \ec{apply}.
\emph{Named identifiers syntactically cited in the proof block:}
\begin{lstlisting}[language=EasyCrypt,basicstyle=\ttfamily\scriptsize]
HAETAE
Pr
ROMInternalTranscriptPaperSimAsNMA
SIG
Samp
UF_NMA
euf_cma_security_from_public_sim_nma_and_counted_bimodal
ler_trans
main
mlwe_bound
msis_bound
paper_nma_hop
public_sim_to_paper_sim_nma_hop_from_sampling_hop
rejection_sampling_loss_term
sampling_hop
\end{lstlisting}

\end{formaltheorem}

\begin{formaltheorem}[EasyCrypt line 704]
\noindent\textbf{EasyCrypt name.}
\begin{lstlisting}[language=EasyCrypt,basicstyle=\ttfamily\scriptsize]
euf_cma_security_from_rom_paper_sim_nma_and_counted_bimodal
\end{lstlisting}
\noindent\textbf{Checked EasyCrypt statement.}
\begin{lstlisting}[language=EasyCrypt,basicstyle=\ttfamily\scriptsize]
lemma euf_cma_security_from_rom_paper_sim_nma_and_counted_bimodal &m :
  Pr[SIG.UF_NMA(H, HAETAE,
       ROMInternalTranscriptPaperSimAsNMA
         (A, ROMPaperSimSigningSampler(H))).main() @ &m : res] <=
    Pr[MLWE_Game(Bmlwe).main() @ &m : res] +
    Pr[BimodalSelfTargetMSIS_Game(Bbimodal).main() @ &m : res] =>
  Pr[MLWE_Game(Bmlwe).main() @ &m : res] <= mlwe_hardness_term =>
  Pr[ModuleSIS_Game(BimodalMSIS_As_ModuleSIS(Bbimodal)).main()
       @ &m : res] <= module_sis_hardness_term =>
  Pr[SIG.EUF_CMA(H, HAETAE, A).main() @ &m : res] <=
    haetae_euf_counted_rom_bound.
\end{lstlisting}
\noindent\textbf{Formal mathematical reading.}
Mathematical theorem. This theorem has the form \(\Pr[G:E] \le B\) is a game-probability statement.  It is one of the exact hops, bad-event bounds, or final advantage inequalities used in the reduction.

\noindent\textbf{Role in the proof.}
Use in this chapter. It contributes directly to an advantage bound, bad-event bound, or real-arithmetic composition step.

\noindent\textbf{Proof-script interpretation.}
Proof-script reading. The machine proof decomposes the theorem into rewrites, program-logic obligations, relational equivalences, and arithmetic side conditions.
\emph{Main tactics visible in the proof script:} \ec{apply}, \ec{rewrite}.
\emph{Named identifiers syntactically cited in the proof block:}
\begin{lstlisting}[language=EasyCrypt,basicstyle=\ttfamily\scriptsize]
euf_cma_security_from_public_sim_nma_and_counted_bimodal
mlwe_bound
msis_bound
nma_hop
rom_internal_nma_public_sim_rom_paper_sim_exact
\end{lstlisting}

\end{formaltheorem}

\begin{formaltheorem}[EasyCrypt line 725]
\noindent\textbf{EasyCrypt name.}
\begin{lstlisting}[language=EasyCrypt,basicstyle=\ttfamily\scriptsize]
euf_cma_security_from_real_signing_paper_sim_nma_and_counted_bimodal
\end{lstlisting}
\noindent\textbf{Checked EasyCrypt statement.}
\begin{lstlisting}[language=EasyCrypt,basicstyle=\ttfamily\scriptsize]
lemma euf_cma_security_from_real_signing_paper_sim_nma_and_counted_bimodal &m :
  Pr[SIG.UF_NMA(H, HAETAE,
       ROMInternalTranscriptPaperSimAsNMA
         (A, RealSigningPaperSimSampler(H))).main() @ &m : res] <=
    Pr[MLWE_Game(Bmlwe).main() @ &m : res] +
    Pr[BimodalSelfTargetMSIS_Game(Bbimodal).main() @ &m : res] =>
  Pr[MLWE_Game(Bmlwe).main() @ &m : res] <= mlwe_hardness_term =>
  Pr[ModuleSIS_Game(BimodalMSIS_As_ModuleSIS(Bbimodal)).main()
       @ &m : res] <= module_sis_hardness_term =>
  Pr[SIG.EUF_CMA(H, HAETAE, A).main() @ &m : res] <=
    haetae_euf_counted_rom_bound.
\end{lstlisting}
\noindent\textbf{Formal mathematical reading.}
Mathematical theorem. This theorem has the form \(\Pr[G:E] \le B\) is a game-probability statement.  It is one of the exact hops, bad-event bounds, or final advantage inequalities used in the reduction.

\noindent\textbf{Role in the proof.}
Use in this chapter. It contributes directly to an advantage bound, bad-event bound, or real-arithmetic composition step.

\noindent\textbf{Proof-script interpretation.}
Proof-script reading. The machine proof decomposes the theorem into rewrites, program-logic obligations, relational equivalences, and arithmetic side conditions.
\emph{Main tactics visible in the proof script:} \ec{apply}, \ec{rewrite}.
\emph{Named identifiers syntactically cited in the proof block:}
\begin{lstlisting}[language=EasyCrypt,basicstyle=\ttfamily\scriptsize]
euf_cma_security_from_structural_nma_and_counted_bimodal
mlwe_bound
msis_bound
nma_hop
rom_internal_nma_real_signing_paper_sim_exact
\end{lstlisting}

\end{formaltheorem}

\begin{formaltheorem}[EasyCrypt line 746]
\noindent\textbf{EasyCrypt name.}
\begin{lstlisting}[language=EasyCrypt,basicstyle=\ttfamily\scriptsize]
euf_cma_security_from_haetae_rejection_paper_sim_nma_and_counted_bimodal
\end{lstlisting}
\noindent\textbf{Checked EasyCrypt statement.}
\begin{lstlisting}[language=EasyCrypt,basicstyle=\ttfamily\scriptsize]
lemma euf_cma_security_from_haetae_rejection_paper_sim_nma_and_counted_bimodal &m :
  Pr[SIG.UF_NMA(H, HAETAE,
       ROMInternalTranscriptPaperSimAsNMA
         (A, HAETAERejectionPaperSimSampler(H))).main() @ &m : res] <=
    Pr[MLWE_Game(Bmlwe).main() @ &m : res] +
    Pr[BimodalSelfTargetMSIS_Game(Bbimodal).main() @ &m : res] =>
  Pr[MLWE_Game(Bmlwe).main() @ &m : res] <= mlwe_hardness_term =>
  Pr[ModuleSIS_Game(BimodalMSIS_As_ModuleSIS(Bbimodal)).main()
       @ &m : res] <= module_sis_hardness_term =>
  Pr[SIG.EUF_CMA(H, HAETAE, A).main() @ &m : res] <=
    haetae_euf_counted_rom_bound.
\end{lstlisting}
\noindent\textbf{Formal mathematical reading.}
Mathematical theorem. This theorem has the form \(\Pr[G:E] \le B\) is a game-probability statement.  It is one of the exact hops, bad-event bounds, or final advantage inequalities used in the reduction.

\noindent\textbf{Role in the proof.}
Use in this chapter. It contributes directly to an advantage bound, bad-event bound, or real-arithmetic composition step.

\noindent\textbf{Proof-script interpretation.}
Proof-script reading. The machine proof decomposes the theorem into rewrites, program-logic obligations, relational equivalences, and arithmetic side conditions.
\emph{Main tactics visible in the proof script:} \ec{apply}, \ec{rewrite}.
\emph{Named identifiers syntactically cited in the proof block:}
\begin{lstlisting}[language=EasyCrypt,basicstyle=\ttfamily\scriptsize]
euf_cma_security_from_real_signing_paper_sim_nma_and_counted_bimodal
mlwe_bound
msis_bound
nma_hop
rom_internal_nma_real_signing_haetae_rejection_paper_sim_exact
\end{lstlisting}

\end{formaltheorem}

\begin{formaltheorem}[EasyCrypt line 768]
\noindent\textbf{EasyCrypt name.}
\begin{lstlisting}[language=EasyCrypt,basicstyle=\ttfamily\scriptsize]
exact_hyperball_haetae_rejection_nma_bound_from_paper_sim
\end{lstlisting}
\noindent\textbf{Checked EasyCrypt statement.}
\begin{lstlisting}[language=EasyCrypt,basicstyle=\ttfamily\scriptsize]
lemma exact_hyperball_haetae_rejection_nma_bound_from_paper_sim &m :
  Pr[SIG.UF_NMA(H, HAETAE,
       ROMInternalTranscriptPaperSimAsNMA
         (A, ExactHyperballPaperSimSampler(H))).main() @ &m : res] <=
    Pr[MLWE_Game(Bmlwe).main() @ &m : res] +
    Pr[BimodalSelfTargetMSIS_Game(Bbimodal).main() @ &m : res] =>
  Pr[SIG.UF_NMA(H, HAETAE,
       ROMInternalTranscriptPaperSimAsNMA
         (A, ExactHyperballHAETAERejectionPaperSimSampler(H))).main()
       @ &m : res] <=
    Pr[MLWE_Game(Bmlwe).main() @ &m : res] +
    Pr[BimodalSelfTargetMSIS_Game(Bbimodal).main() @ &m : res].
\end{lstlisting}
\noindent\textbf{Formal mathematical reading.}
Mathematical theorem. This theorem has the form \(\Pr[G:E] \le B\) is a game-probability statement.  It is one of the exact hops, bad-event bounds, or final advantage inequalities used in the reduction.

\noindent\textbf{Role in the proof.}
Use in this chapter. It contributes directly to an advantage bound, bad-event bound, or real-arithmetic composition step.

\noindent\textbf{Proof-script interpretation.}
Proof-script reading. The machine proof decomposes the theorem into rewrites, program-logic obligations, relational equivalences, and arithmetic side conditions.
\emph{Main tactics visible in the proof script:} \ec{rewrite}, \ec{apply}.
\emph{Named identifiers syntactically cited in the proof block:}
\begin{lstlisting}[language=EasyCrypt,basicstyle=\ttfamily\scriptsize]
nma_hop
rom_internal_nma_exact_hyperball_rejection_paper_sim_no_fallback_exact
\end{lstlisting}

\end{formaltheorem}

\begin{formaltheorem}[EasyCrypt line 787]
\noindent\textbf{EasyCrypt name.}
\begin{lstlisting}[language=EasyCrypt,basicstyle=\ttfamily\scriptsize]
haetae_rejection_nma_bound_from_exact_hyperball_sampling_hop
\end{lstlisting}
\noindent\textbf{Checked EasyCrypt statement.}
\begin{lstlisting}[language=EasyCrypt,basicstyle=\ttfamily\scriptsize]
lemma haetae_rejection_nma_bound_from_exact_hyperball_sampling_hop &m :
  Pr[SIG.UF_NMA(H, HAETAE,
       ROMInternalTranscriptPaperSimAsNMA
         (A, HAETAERejectionPaperSimSampler(H))).main() @ &m : res] <=
  Pr[SIG.UF_NMA(H, HAETAE,
       ROMInternalTranscriptPaperSimAsNMA
         (A, ExactHyperballHAETAERejectionPaperSimSampler(H))).main()
       @ &m : res] + rejection_sampling_loss_term =>
  Pr[SIG.UF_NMA(H, HAETAE,
       ROMInternalTranscriptPaperSimAsNMA
         (A, ExactHyperballPaperSimSampler(H))).main() @ &m : res] +
    rejection_sampling_loss_term <=
    Pr[MLWE_Game(Bmlwe).main() @ &m : res] +
    Pr[BimodalSelfTargetMSIS_Game(Bbimodal).main() @ &m : res] =>
  Pr[SIG.UF_NMA(H, HAETAE,
       ROMInternalTranscriptPaperSimAsNMA
         (A, HAETAERejectionPaperSimSampler(H))).main() @ &m : res] <=
    Pr[MLWE_Game(Bmlwe).main() @ &m : res] +
    Pr[BimodalSelfTargetMSIS_Game(Bbimodal).main() @ &m : res].
\end{lstlisting}
\noindent\textbf{Formal mathematical reading.}
Mathematical theorem. This theorem has the form \(\Pr[G:E] \le B\) is a game-probability statement.  It is one of the exact hops, bad-event bounds, or final advantage inequalities used in the reduction.

\noindent\textbf{Role in the proof.}
Use in this chapter. It contributes directly to an advantage bound, bad-event bound, or real-arithmetic composition step.

\noindent\textbf{Proof-script interpretation.}
Proof-script reading. The machine proof decomposes the theorem into rewrites, program-logic obligations, relational equivalences, and arithmetic side conditions.
\emph{Main tactics visible in the proof script:} \ec{apply}, \ec{rewrite}.
\emph{Named identifiers syntactically cited in the proof block:}
\begin{lstlisting}[language=EasyCrypt,basicstyle=\ttfamily\scriptsize]
ExactHyperballHAETAERejectionPaperSimSampler
HAETAE
Pr
ROMInternalTranscriptPaperSimAsNMA
SIG
UF_NMA
exact_nma_hop
ler_trans
main
rejection_sampling_loss_term
rom_internal_nma_exact_hyperball_rejection_paper_sim_no_fallback_exact
sampling_hop
\end{lstlisting}

\end{formaltheorem}

\begin{formaltheorem}[EasyCrypt line 819]
\noindent\textbf{EasyCrypt name.}
\begin{lstlisting}[language=EasyCrypt,basicstyle=\ttfamily\scriptsize]
haetae_rejection_exact_hyperball_sampling_hop_from_real_signing_exact_hyperball_hop
\end{lstlisting}
\noindent\textbf{Checked EasyCrypt statement.}
\begin{lstlisting}[language=EasyCrypt,basicstyle=\ttfamily\scriptsize]
lemma haetae_rejection_exact_hyperball_sampling_hop_from_real_signing_exact_hyperball_hop &m :
  Pr[SIG.UF_NMA(H, HAETAE,
       ROMInternalTranscriptPaperSimAsNMA
         (A, RealSigningPaperSimSampler(H))).main() @ &m : res] <=
  Pr[SIG.UF_NMA(H, HAETAE,
       ROMInternalTranscriptPaperSimAsNMA
         (A, ExactHyperballPaperSimSampler(H))).main() @ &m : res] +
    rejection_sampling_loss_term =>
  Pr[SIG.UF_NMA(H, HAETAE,
       ROMInternalTranscriptPaperSimAsNMA
         (A, HAETAERejectionPaperSimSampler(H))).main() @ &m : res] <=
  Pr[SIG.UF_NMA(H, HAETAE,
       ROMInternalTranscriptPaperSimAsNMA
         (A, ExactHyperballHAETAERejectionPaperSimSampler(H))).main()
       @ &m : res] + rejection_sampling_loss_term.
\end{lstlisting}
\noindent\textbf{Formal mathematical reading.}
Mathematical theorem. This theorem has the form \(\Pr[G:E] \le B\) is a game-probability statement.  It is one of the exact hops, bad-event bounds, or final advantage inequalities used in the reduction.

\noindent\textbf{Role in the proof.}
Use in this chapter. It proves an exact game replacement or simulator equivalence.

\noindent\textbf{Proof-script interpretation.}
Proof-script reading. The machine proof decomposes the theorem into rewrites, program-logic obligations, relational equivalences, and arithmetic side conditions.
\emph{Main tactics visible in the proof script:} \ec{rewrite}, \ec{apply}.
\emph{Named identifiers syntactically cited in the proof block:}
\begin{lstlisting}[language=EasyCrypt,basicstyle=\ttfamily\scriptsize]
real_exact_hop
rom_internal_nma_exact_hyperball_rejection_paper_sim_no_fallback_exact
rom_internal_nma_real_signing_haetae_rejection_paper_sim_exact
\end{lstlisting}

\end{formaltheorem}

\begin{formaltheorem}[EasyCrypt line 843]
\noindent\textbf{EasyCrypt name.}
\begin{lstlisting}[language=EasyCrypt,basicstyle=\ttfamily\scriptsize]
haetae_rejection_ro_exact_hyperball_sampling_hop_from_real_signing_ro_exact_hyperball_hop
\end{lstlisting}
\noindent\textbf{Checked EasyCrypt statement.}
\begin{lstlisting}[language=EasyCrypt,basicstyle=\ttfamily\scriptsize]
lemma haetae_rejection_ro_exact_hyperball_sampling_hop_from_real_signing_ro_exact_hyperball_hop &m :
  Pr[SIG.UF_NMA(H, HAETAE,
       ROMInternalTranscriptPaperSimAsNMA
         (A, RealSigningPaperSimSampler(H))).main() @ &m : res] <=
  Pr[SIG.UF_NMA(H, HAETAE,
       ROMInternalTranscriptPaperSimAsNMA
         (A, ROExactHyperballPaperSimSampler(H))).main() @ &m : res] +
    rejection_sampling_loss_term =>
  Pr[SIG.UF_NMA(H, HAETAE,
       ROMInternalTranscriptPaperSimAsNMA
         (A, HAETAERejectionPaperSimSampler(H))).main() @ &m : res] <=
  Pr[SIG.UF_NMA(H, HAETAE,
       ROMInternalTranscriptPaperSimAsNMA
         (A, ROExactHyperballPaperSimSampler(H))).main() @ &m : res] +
    rejection_sampling_loss_term.
\end{lstlisting}
\noindent\textbf{Formal mathematical reading.}
Mathematical theorem. This theorem has the form \(\Pr[G:E] \le B\) is a game-probability statement.  It is one of the exact hops, bad-event bounds, or final advantage inequalities used in the reduction.

\noindent\textbf{Role in the proof.}
Use in this chapter. It proves an exact game replacement or simulator equivalence.

\noindent\textbf{Proof-script interpretation.}
Proof-script reading. The machine proof decomposes the theorem into rewrites, program-logic obligations, relational equivalences, and arithmetic side conditions.
\emph{Main tactics visible in the proof script:} \ec{rewrite}, \ec{apply}.
\emph{Named identifiers syntactically cited in the proof block:}
\begin{lstlisting}[language=EasyCrypt,basicstyle=\ttfamily\scriptsize]
real_ro_exact_hop
rom_internal_nma_real_signing_haetae_rejection_paper_sim_exact
\end{lstlisting}

\end{formaltheorem}

\begin{formaltheorem}[EasyCrypt line 865]
\noindent\textbf{EasyCrypt name.}
\begin{lstlisting}[language=EasyCrypt,basicstyle=\ttfamily\scriptsize]
real_signing_ro_exact_hyperball_hop_from_ro_signing_attempt_hop
\end{lstlisting}
\noindent\textbf{Checked EasyCrypt statement.}
\begin{lstlisting}[language=EasyCrypt,basicstyle=\ttfamily\scriptsize]
lemma real_signing_ro_exact_hyperball_hop_from_ro_signing_attempt_hop &m :
  Pr[SIG.UF_NMA(H, HAETAE,
       ROMInternalTranscriptPaperSimAsNMA
         (A, ROSigningAttemptPaperSimSampler(H))).main() @ &m : res] <=
  Pr[SIG.UF_NMA(H, HAETAE,
       ROMInternalTranscriptPaperSimAsNMA
         (A, ROExactHyperballPaperSimSampler(H))).main() @ &m : res] +
    rejection_sampling_loss_term =>
  Pr[SIG.UF_NMA(H, HAETAE,
       ROMInternalTranscriptPaperSimAsNMA
         (A, RealSigningPaperSimSampler(H))).main() @ &m : res] <=
  Pr[SIG.UF_NMA(H, HAETAE,
       ROMInternalTranscriptPaperSimAsNMA
         (A, ROExactHyperballPaperSimSampler(H))).main() @ &m : res] +
    rejection_sampling_loss_term.
\end{lstlisting}
\noindent\textbf{Formal mathematical reading.}
Mathematical theorem. This theorem has the form \(\Pr[G:E] \le B\) is a game-probability statement.  It is one of the exact hops, bad-event bounds, or final advantage inequalities used in the reduction.

\noindent\textbf{Role in the proof.}
Use in this chapter. It proves an exact game replacement or simulator equivalence.

\noindent\textbf{Proof-script interpretation.}
Proof-script reading. The machine proof decomposes the theorem into rewrites, program-logic obligations, relational equivalences, and arithmetic side conditions.
\emph{Main tactics visible in the proof script:} \ec{rewrite}, \ec{apply}.
\emph{Named identifiers syntactically cited in the proof block:}
\begin{lstlisting}[language=EasyCrypt,basicstyle=\ttfamily\scriptsize]
ro_attempt_hop
rom_internal_nma_real_signing_ro_attempt_paper_sim_exact
\end{lstlisting}

\end{formaltheorem}

\begin{formaltheorem}[EasyCrypt line 886]
\noindent\textbf{EasyCrypt name.}
\begin{lstlisting}[language=EasyCrypt,basicstyle=\ttfamily\scriptsize]
real_signing_ro_exact_hyperball_hop_from_ro_signing_attempt_hop_loss
\end{lstlisting}
\noindent\textbf{Checked EasyCrypt statement.}
\begin{lstlisting}[language=EasyCrypt,basicstyle=\ttfamily\scriptsize]
lemma real_signing_ro_exact_hyperball_hop_from_ro_signing_attempt_hop_loss
  &m (loss : real) :
  Pr[SIG.UF_NMA(H, HAETAE,
       ROMInternalTranscriptPaperSimAsNMA
         (A, ROSigningAttemptPaperSimSampler(H))).main() @ &m : res] <=
  Pr[SIG.UF_NMA(H, HAETAE,
       ROMInternalTranscriptPaperSimAsNMA
         (A, ROExactHyperballPaperSimSampler(H))).main() @ &m : res] +
    loss =>
  Pr[SIG.UF_NMA(H, HAETAE,
       ROMInternalTranscriptPaperSimAsNMA
         (A, RealSigningPaperSimSampler(H))).main() @ &m : res] <=
  Pr[SIG.UF_NMA(H, HAETAE,
       ROMInternalTranscriptPaperSimAsNMA
         (A, ROExactHyperballPaperSimSampler(H))).main() @ &m : res] +
    loss.
\end{lstlisting}
\noindent\textbf{Formal mathematical reading.}
Mathematical theorem. This theorem has the form \(\Pr[G:E] \le B\) is a game-probability statement.  It is one of the exact hops, bad-event bounds, or final advantage inequalities used in the reduction.

\noindent\textbf{Role in the proof.}
Use in this chapter. It contributes directly to an advantage bound, bad-event bound, or real-arithmetic composition step.

\noindent\textbf{Proof-script interpretation.}
Proof-script reading. The machine proof decomposes the theorem into rewrites, program-logic obligations, relational equivalences, and arithmetic side conditions.
\emph{Main tactics visible in the proof script:} \ec{rewrite}, \ec{apply}.
\emph{Named identifiers syntactically cited in the proof block:}
\begin{lstlisting}[language=EasyCrypt,basicstyle=\ttfamily\scriptsize]
ro_attempt_hop
rom_internal_nma_real_signing_ro_attempt_paper_sim_exact
\end{lstlisting}

\end{formaltheorem}

\begin{formaltheorem}[EasyCrypt line 908]
\noindent\textbf{EasyCrypt name.}
\begin{lstlisting}[language=EasyCrypt,basicstyle=\ttfamily\scriptsize]
real_signing_ro_exact_hyperball_hop_from_budgeted_lifting
\end{lstlisting}
\noindent\textbf{Checked EasyCrypt statement.}
\begin{lstlisting}[language=EasyCrypt,basicstyle=\ttfamily\scriptsize]
lemma real_signing_ro_exact_hyperball_hop_from_budgeted_lifting &m :
  Pr[SIG.UF_NMA(H, HAETAE,
       ROMInternalTranscriptPaperSimAsNMA
         (A, ROSigningAttemptPaperSimSampler(H))).main() @ &m : res] <=
  Pr[SIG.UF_NMA(H, HAETAE,
       ROMInternalTranscriptBudgetedPaperSimAsNMA
         (A, ROSigningAttemptPaperSimSampler(H))).main() @ &m : res] =>
  Pr[SIG.UF_NMA(H, HAETAE,
       ROMInternalTranscriptBudgetedPaperSimAsNMA
         (A, ROExactHyperballPaperSimSampler(H))).main() @ &m : res] <=
  Pr[SIG.UF_NMA(H, HAETAE,
       ROMInternalTranscriptPaperSimAsNMA
         (A, ROExactHyperballPaperSimSampler(H))).main() @ &m : res] =>
  structural_to_exact_hyperball_paper_sample_loss_obligation haetae_mode =>
  Pr[SIG.UF_NMA(H, HAETAE,
       ROMInternalTranscriptPaperSimAsNMA
         (A, RealSigningPaperSimSampler(H))).main() @ &m : res] <=
  Pr[SIG.UF_NMA(H, HAETAE,
       ROMInternalTranscriptPaperSimAsNMA
         (A, ROExactHyperballPaperSimSampler(H))).main() @ &m : res] +
    rom_signature_query_budget * rejection_sampling_loss_term.
\end{lstlisting}
\noindent\textbf{Formal mathematical reading.}
Mathematical theorem. This theorem has the form \(\Pr[G:E] \le B\) is a game-probability statement.  It is one of the exact hops, bad-event bounds, or final advantage inequalities used in the reduction.

\noindent\textbf{Role in the proof.}
Use in this chapter. It proves an exact game replacement or simulator equivalence.

\noindent\textbf{Proof-script interpretation.}
Proof-script reading. The machine proof decomposes the theorem into rewrites, program-logic obligations, relational equivalences, and arithmetic side conditions.
\emph{Main tactics visible in the proof script:} \ec{apply}.
\emph{Named identifiers syntactically cited in the proof block:}
\begin{lstlisting}[language=EasyCrypt,basicstyle=\ttfamily\scriptsize]
attempt_unbudgeted_to_budgeted
exact_budgeted_to_unbudgeted
real_signing_ro_exact_hyperball_hop_from_ro_signing_attempt_hop_loss
rejection_sampling_loss_term
rom_internal_nma_ro_signing_attempt_ro_exact_hyperball_loss_from_budgeted_lifting
rom_signature_query_budget
sample_loss
\end{lstlisting}

\end{formaltheorem}

\begin{formaltheorem}[EasyCrypt line 943]
\noindent\textbf{EasyCrypt name.}
\begin{lstlisting}[language=EasyCrypt,basicstyle=\ttfamily\scriptsize]
haetae_rejection_ro_exact_hyperball_sampling_hop_from_budgeted_lifting
\end{lstlisting}
\noindent\textbf{Checked EasyCrypt statement.}
\begin{lstlisting}[language=EasyCrypt,basicstyle=\ttfamily\scriptsize]
lemma haetae_rejection_ro_exact_hyperball_sampling_hop_from_budgeted_lifting
  &m :
  Pr[SIG.UF_NMA(H, HAETAE,
       ROMInternalTranscriptPaperSimAsNMA
         (A, ROSigningAttemptPaperSimSampler(H))).main() @ &m : res] <=
  Pr[SIG.UF_NMA(H, HAETAE,
       ROMInternalTranscriptBudgetedPaperSimAsNMA
         (A, ROSigningAttemptPaperSimSampler(H))).main() @ &m : res] =>
  Pr[SIG.UF_NMA(H, HAETAE,
       ROMInternalTranscriptBudgetedPaperSimAsNMA
         (A, ROExactHyperballPaperSimSampler(H))).main() @ &m : res] <=
  Pr[SIG.UF_NMA(H, HAETAE,
       ROMInternalTranscriptPaperSimAsNMA
         (A, ROExactHyperballPaperSimSampler(H))).main() @ &m : res] =>
  structural_to_exact_hyperball_paper_sample_loss_obligation haetae_mode =>
  Pr[SIG.UF_NMA(H, HAETAE,
       ROMInternalTranscriptPaperSimAsNMA
         (A, HAETAERejectionPaperSimSampler(H))).main() @ &m : res] <=
  Pr[SIG.UF_NMA(H, HAETAE,
       ROMInternalTranscriptPaperSimAsNMA
         (A, ROExactHyperballPaperSimSampler(H))).main() @ &m : res] +
    rom_signature_query_budget * rejection_sampling_loss_term.
\end{lstlisting}
\noindent\textbf{Formal mathematical reading.}
Mathematical theorem. This theorem has the form \(\Pr[G:E] \le B\) is a game-probability statement.  It is one of the exact hops, bad-event bounds, or final advantage inequalities used in the reduction.

\noindent\textbf{Role in the proof.}
Use in this chapter. It proves an exact game replacement or simulator equivalence.

\noindent\textbf{Proof-script interpretation.}
Proof-script reading. The machine proof decomposes the theorem into rewrites, program-logic obligations, relational equivalences, and arithmetic side conditions.
\emph{Main tactics visible in the proof script:} \ec{rewrite}, \ec{apply}.
\emph{Named identifiers syntactically cited in the proof block:}
\begin{lstlisting}[language=EasyCrypt,basicstyle=\ttfamily\scriptsize]
attempt_unbudgeted_to_budgeted
exact_budgeted_to_unbudgeted
real_signing_ro_exact_hyperball_hop_from_budgeted_lifting
rom_internal_nma_real_signing_haetae_rejection_paper_sim_exact
sample_loss
\end{lstlisting}

\end{formaltheorem}

\begin{formaltheorem}[EasyCrypt line 974]
\noindent\textbf{EasyCrypt name.}
\begin{lstlisting}[language=EasyCrypt,basicstyle=\ttfamily\scriptsize]
euf_cma_security_from_exact_hyperball_paper_sim_bridge_and_counted_bimodal
\end{lstlisting}
\noindent\textbf{Checked EasyCrypt statement.}
\begin{lstlisting}[language=EasyCrypt,basicstyle=\ttfamily\scriptsize]
lemma euf_cma_security_from_exact_hyperball_paper_sim_bridge_and_counted_bimodal &m :
  Pr[SIG.UF_NMA(H, HAETAE,
       ROMInternalTranscriptPaperSimAsNMA
         (A, HAETAERejectionPaperSimSampler(H))).main() @ &m : res] <=
  Pr[SIG.UF_NMA(H, HAETAE,
       ROMInternalTranscriptPaperSimAsNMA
         (A, ExactHyperballHAETAERejectionPaperSimSampler(H))).main()
       @ &m : res] =>
  Pr[SIG.UF_NMA(H, HAETAE,
       ROMInternalTranscriptPaperSimAsNMA
         (A, ExactHyperballPaperSimSampler(H))).main() @ &m : res] <=
    Pr[MLWE_Game(Bmlwe).main() @ &m : res] +
    Pr[BimodalSelfTargetMSIS_Game(Bbimodal).main() @ &m : res] =>
  Pr[MLWE_Game(Bmlwe).main() @ &m : res] <= mlwe_hardness_term =>
  Pr[ModuleSIS_Game(BimodalMSIS_As_ModuleSIS(Bbimodal)).main()
       @ &m : res] <= module_sis_hardness_term =>
  Pr[SIG.EUF_CMA(H, HAETAE, A).main() @ &m : res] <=
    haetae_euf_counted_rom_bound.
\end{lstlisting}
\noindent\textbf{Formal mathematical reading.}
Mathematical theorem. This theorem has the form \(\Pr[G:E] \le B\) is a game-probability statement.  It is one of the exact hops, bad-event bounds, or final advantage inequalities used in the reduction.

\noindent\textbf{Role in the proof.}
Use in this chapter. It contributes directly to an advantage bound, bad-event bound, or real-arithmetic composition step.

\noindent\textbf{Proof-script interpretation.}
Proof-script reading. The machine proof decomposes the theorem into rewrites, program-logic obligations, relational equivalences, and arithmetic side conditions.
\emph{Main tactics visible in the proof script:} \ec{apply}.
\emph{Named identifiers syntactically cited in the proof block:}
\begin{lstlisting}[language=EasyCrypt,basicstyle=\ttfamily\scriptsize]
ExactHyperballHAETAERejectionPaperSimSampler
HAETAE
Pr
ROMInternalTranscriptPaperSimAsNMA
SIG
UF_NMA
euf_cma_security_from_haetae_rejection_paper_sim_nma_and_counted_bimodal
exact_bridge
exact_hyperball_haetae_rejection_nma_bound_from_paper_sim
exact_nma_hop
ler_trans
main
mlwe_bound
msis_bound
\end{lstlisting}

\end{formaltheorem}

\begin{formaltheorem}[EasyCrypt line 1007]
\noindent\textbf{EasyCrypt name.}
\begin{lstlisting}[language=EasyCrypt,basicstyle=\ttfamily\scriptsize]
euf_cma_security_from_exact_hyperball_sampling_hop_and_counted_bimodal
\end{lstlisting}
\noindent\textbf{Checked EasyCrypt statement.}
\begin{lstlisting}[language=EasyCrypt,basicstyle=\ttfamily\scriptsize]
lemma euf_cma_security_from_exact_hyperball_sampling_hop_and_counted_bimodal &m :
  Pr[SIG.UF_NMA(H, HAETAE,
       ROMInternalTranscriptPaperSimAsNMA
         (A, HAETAERejectionPaperSimSampler(H))).main() @ &m : res] <=
  Pr[SIG.UF_NMA(H, HAETAE,
       ROMInternalTranscriptPaperSimAsNMA
         (A, ExactHyperballHAETAERejectionPaperSimSampler(H))).main()
       @ &m : res] + rejection_sampling_loss_term =>
  Pr[SIG.UF_NMA(H, HAETAE,
       ROMInternalTranscriptPaperSimAsNMA
         (A, ExactHyperballPaperSimSampler(H))).main() @ &m : res] +
    rejection_sampling_loss_term <=
    Pr[MLWE_Game(Bmlwe).main() @ &m : res] +
    Pr[BimodalSelfTargetMSIS_Game(Bbimodal).main() @ &m : res] =>
  Pr[MLWE_Game(Bmlwe).main() @ &m : res] <= mlwe_hardness_term =>
  Pr[ModuleSIS_Game(BimodalMSIS_As_ModuleSIS(Bbimodal)).main()
       @ &m : res] <= module_sis_hardness_term =>
  Pr[SIG.EUF_CMA(H, HAETAE, A).main() @ &m : res] <=
    haetae_euf_counted_rom_bound.
\end{lstlisting}
\noindent\textbf{Formal mathematical reading.}
Mathematical theorem. This theorem has the form \(\Pr[G:E] \le B\) is a game-probability statement.  It is one of the exact hops, bad-event bounds, or final advantage inequalities used in the reduction.

\noindent\textbf{Role in the proof.}
Use in this chapter. It contributes directly to an advantage bound, bad-event bound, or real-arithmetic composition step.

\noindent\textbf{Proof-script interpretation.}
Proof-script reading. The machine proof decomposes the theorem into rewrites, program-logic obligations, relational equivalences, and arithmetic side conditions.
\emph{Main tactics visible in the proof script:} \ec{apply}.
\emph{Named identifiers syntactically cited in the proof block:}
\begin{lstlisting}[language=EasyCrypt,basicstyle=\ttfamily\scriptsize]
euf_cma_security_from_haetae_rejection_paper_sim_nma_and_counted_bimodal
exact_nma_hop
haetae_rejection_nma_bound_from_exact_hyperball_sampling_hop
mlwe_bound
msis_bound
sampling_hop
\end{lstlisting}

\end{formaltheorem}

\begin{formaltheorem}[EasyCrypt line 1036]
\noindent\textbf{EasyCrypt name.}
\begin{lstlisting}[language=EasyCrypt,basicstyle=\ttfamily\scriptsize]
euf_cma_security_from_ro_exact_hyperball_sampling_hop_and_counted_bimodal
\end{lstlisting}
\noindent\textbf{Checked EasyCrypt statement.}
\begin{lstlisting}[language=EasyCrypt,basicstyle=\ttfamily\scriptsize]
lemma euf_cma_security_from_ro_exact_hyperball_sampling_hop_and_counted_bimodal &m :
  Pr[SIG.UF_NMA(H, HAETAE,
       ROMInternalTranscriptPaperSimAsNMA
         (A, RealSigningPaperSimSampler(H))).main() @ &m : res] <=
  Pr[SIG.UF_NMA(H, HAETAE,
       ROMInternalTranscriptPaperSimAsNMA
         (A, ROExactHyperballPaperSimSampler(H))).main() @ &m : res] +
    rejection_sampling_loss_term =>
  Pr[SIG.UF_NMA(H, HAETAE,
       ROMInternalTranscriptPaperSimAsNMA
         (A, ROExactHyperballPaperSimSampler(H))).main() @ &m : res] +
    rejection_sampling_loss_term <=
    Pr[MLWE_Game(Bmlwe).main() @ &m : res] +
    Pr[BimodalSelfTargetMSIS_Game(Bbimodal).main() @ &m : res] =>
  Pr[MLWE_Game(Bmlwe).main() @ &m : res] <= mlwe_hardness_term =>
  Pr[ModuleSIS_Game(BimodalMSIS_As_ModuleSIS(Bbimodal)).main()
       @ &m : res] <= module_sis_hardness_term =>
  Pr[SIG.EUF_CMA(H, HAETAE, A).main() @ &m : res] <=
    haetae_euf_counted_rom_bound.
\end{lstlisting}
\noindent\textbf{Formal mathematical reading.}
Mathematical theorem. This theorem has the form \(\Pr[G:E] \le B\) is a game-probability statement.  It is one of the exact hops, bad-event bounds, or final advantage inequalities used in the reduction.

\noindent\textbf{Role in the proof.}
Use in this chapter. It contributes directly to an advantage bound, bad-event bound, or real-arithmetic composition step.

\noindent\textbf{Proof-script interpretation.}
Proof-script reading. The machine proof decomposes the theorem into rewrites, program-logic obligations, relational equivalences, and arithmetic side conditions.
\emph{Main tactics visible in the proof script:} \ec{apply}.
\emph{Named identifiers syntactically cited in the proof block:}
\begin{lstlisting}[language=EasyCrypt,basicstyle=\ttfamily\scriptsize]
HAETAE
Pr
ROExactHyperballPaperSimSampler
ROMInternalTranscriptPaperSimAsNMA
SIG
UF_NMA
euf_cma_security_from_haetae_rejection_paper_sim_nma_and_counted_bimodal
haetae_rejection_ro_exact_hyperball_sampling_hop_from_real_signing_ro_exact_hyperball_hop
ler_trans
main
mlwe_bound
msis_bound
real_ro_exact_hop
rejection_sampling_loss_term
ro_exact_nma_hop
\end{lstlisting}

\end{formaltheorem}

\begin{formaltheorem}[EasyCrypt line 1072]
\noindent\textbf{EasyCrypt name.}
\begin{lstlisting}[language=EasyCrypt,basicstyle=\ttfamily\scriptsize]
euf_cma_security_from_ro_signing_attempt_to_ro_exact_hyperball_hop_and_counted_bimodal
\end{lstlisting}
\noindent\textbf{Checked EasyCrypt statement.}
\begin{lstlisting}[language=EasyCrypt,basicstyle=\ttfamily\scriptsize]
lemma euf_cma_security_from_ro_signing_attempt_to_ro_exact_hyperball_hop_and_counted_bimodal &m :
  Pr[SIG.UF_NMA(H, HAETAE,
       ROMInternalTranscriptPaperSimAsNMA
         (A, ROSigningAttemptPaperSimSampler(H))).main() @ &m : res] <=
  Pr[SIG.UF_NMA(H, HAETAE,
       ROMInternalTranscriptPaperSimAsNMA
         (A, ROExactHyperballPaperSimSampler(H))).main() @ &m : res] +
    rejection_sampling_loss_term =>
  Pr[SIG.UF_NMA(H, HAETAE,
       ROMInternalTranscriptPaperSimAsNMA
         (A, ROExactHyperballPaperSimSampler(H))).main() @ &m : res] +
    rejection_sampling_loss_term <=
    Pr[MLWE_Game(Bmlwe).main() @ &m : res] +
    Pr[BimodalSelfTargetMSIS_Game(Bbimodal).main() @ &m : res] =>
  Pr[MLWE_Game(Bmlwe).main() @ &m : res] <= mlwe_hardness_term =>
  Pr[ModuleSIS_Game(BimodalMSIS_As_ModuleSIS(Bbimodal)).main()
       @ &m : res] <= module_sis_hardness_term =>
  Pr[SIG.EUF_CMA(H, HAETAE, A).main() @ &m : res] <=
    haetae_euf_counted_rom_bound.
\end{lstlisting}
\noindent\textbf{Formal mathematical reading.}
Mathematical theorem. This theorem has the form \(\Pr[G:E] \le B\) is a game-probability statement.  It is one of the exact hops, bad-event bounds, or final advantage inequalities used in the reduction.

\noindent\textbf{Role in the proof.}
Use in this chapter. It contributes directly to an advantage bound, bad-event bound, or real-arithmetic composition step.

\noindent\textbf{Proof-script interpretation.}
Proof-script reading. The machine proof decomposes the theorem into rewrites, program-logic obligations, relational equivalences, and arithmetic side conditions.
\emph{Main tactics visible in the proof script:} \ec{apply}.
\emph{Named identifiers syntactically cited in the proof block:}
\begin{lstlisting}[language=EasyCrypt,basicstyle=\ttfamily\scriptsize]
euf_cma_security_from_ro_exact_hyperball_sampling_hop_and_counted_bimodal
mlwe_bound
msis_bound
real_signing_ro_exact_hyperball_hop_from_ro_signing_attempt_hop
ro_attempt_hop
ro_exact_nma_hop
\end{lstlisting}

\end{formaltheorem}

\begin{formaltheorem}[EasyCrypt line 1102]
\noindent\textbf{EasyCrypt name.}
\begin{lstlisting}[language=EasyCrypt,basicstyle=\ttfamily\scriptsize]
real_signing_ro_exact_hyperball_hop_checked_from_loss_budget
\end{lstlisting}
\noindent\textbf{Checked EasyCrypt statement.}
\begin{lstlisting}[language=EasyCrypt,basicstyle=\ttfamily\scriptsize]
lemma real_signing_ro_exact_hyperball_hop_checked_from_loss_budget &m :
  Pr[SIG.UF_NMA(H, HAETAE,
       ROMInternalTranscriptPaperSimAsNMA
         (A, RealSigningPaperSimSampler(H))).main() @ &m : res] <=
  Pr[SIG.UF_NMA(H, HAETAE,
       ROMInternalTranscriptPaperSimAsNMA
         (A, ROExactHyperballPaperSimSampler(H))).main() @ &m : res] +
    rejection_sampling_loss_term.
\end{lstlisting}
\noindent\textbf{Formal mathematical reading.}
Mathematical theorem. This theorem has the form \(\Pr[G:E] \le B\) is a game-probability statement.  It is one of the exact hops, bad-event bounds, or final advantage inequalities used in the reduction.

\noindent\textbf{Role in the proof.}
Use in this chapter. It contributes directly to an advantage bound, bad-event bound, or real-arithmetic composition step.

\noindent\textbf{Proof-script interpretation.}
Proof-script reading. The machine proof decomposes the theorem into rewrites, program-logic obligations, relational equivalences, and arithmetic side conditions.
\emph{Main tactics visible in the proof script:} \ec{apply}.
\emph{Named identifiers syntactically cited in the proof block:}
\begin{lstlisting}[language=EasyCrypt,basicstyle=\ttfamily\scriptsize]
real_signing_ro_exact_hyperball_hop_from_ro_signing_attempt_hop
rom_internal_nma_ro_signing_attempt_ro_exact_hyperball_loss_bound
\end{lstlisting}

\end{formaltheorem}

\begin{formaltheorem}[EasyCrypt line 1116]
\noindent\textbf{EasyCrypt name.}
\begin{lstlisting}[language=EasyCrypt,basicstyle=\ttfamily\scriptsize]
real_signing_ro_exact_hyperball_hop_from_paper_sample_lifting
\end{lstlisting}
\noindent\textbf{Checked EasyCrypt statement.}
\begin{lstlisting}[language=EasyCrypt,basicstyle=\ttfamily\scriptsize]
lemma real_signing_ro_exact_hyperball_hop_from_paper_sample_lifting &m :
  structural_to_exact_hyperball_paper_sample_loss_obligation haetae_mode =>
  Pr[SIG.UF_NMA(H, HAETAE,
       ROMInternalTranscriptPaperSimAsNMA
         (A, ROSigningAttemptPaperSimSampler(H))).main() @ &m : res] <=
  Pr[SIG.UF_NMA(H, HAETAE,
       ROMInternalTranscriptPaperSimAsNMA
         (A, ROExactHyperballPaperSimSampler(H))).main() @ &m : res] +
    rejection_sampling_loss_term =>
  Pr[SIG.UF_NMA(H, HAETAE,
       ROMInternalTranscriptPaperSimAsNMA
         (A, RealSigningPaperSimSampler(H))).main() @ &m : res] <=
  Pr[SIG.UF_NMA(H, HAETAE,
       ROMInternalTranscriptPaperSimAsNMA
         (A, ROExactHyperballPaperSimSampler(H))).main() @ &m : res] +
    rejection_sampling_loss_term.
\end{lstlisting}
\noindent\textbf{Formal mathematical reading.}
Mathematical theorem. This theorem has the form \(\Pr[G:E] \le B\) is a game-probability statement.  It is one of the exact hops, bad-event bounds, or final advantage inequalities used in the reduction.

\noindent\textbf{Role in the proof.}
Use in this chapter. It contributes directly to an advantage bound, bad-event bound, or real-arithmetic composition step.

\noindent\textbf{Proof-script interpretation.}
Proof-script reading. The machine proof decomposes the theorem into rewrites, program-logic obligations, relational equivalences, and arithmetic side conditions.
\emph{Main tactics visible in the proof script:} \ec{apply}.
\emph{Named identifiers syntactically cited in the proof block:}
\begin{lstlisting}[language=EasyCrypt,basicstyle=\ttfamily\scriptsize]
lifted_loss
real_signing_ro_exact_hyperball_hop_from_ro_signing_attempt_hop
rom_internal_nma_ro_signing_attempt_ro_exact_hyperball_loss_from_paper_sample_lifting
sample_loss
\end{lstlisting}

\end{formaltheorem}

\begin{formaltheorem}[EasyCrypt line 1140]
\noindent\textbf{EasyCrypt name.}
\begin{lstlisting}[language=EasyCrypt,basicstyle=\ttfamily\scriptsize]
euf_cma_security_from_paper_sample_lifting_and_counted_bimodal
\end{lstlisting}
\noindent\textbf{Checked EasyCrypt statement.}
\begin{lstlisting}[language=EasyCrypt,basicstyle=\ttfamily\scriptsize]
lemma euf_cma_security_from_paper_sample_lifting_and_counted_bimodal &m :
  structural_to_exact_hyperball_paper_sample_loss_obligation haetae_mode =>
  Pr[SIG.UF_NMA(H, HAETAE,
       ROMInternalTranscriptPaperSimAsNMA
         (A, ROSigningAttemptPaperSimSampler(H))).main() @ &m : res] <=
  Pr[SIG.UF_NMA(H, HAETAE,
       ROMInternalTranscriptPaperSimAsNMA
         (A, ROExactHyperballPaperSimSampler(H))).main() @ &m : res] +
    rejection_sampling_loss_term =>
  Pr[SIG.UF_NMA(H, HAETAE,
       ROMInternalTranscriptPaperSimAsNMA
         (A, ROExactHyperballPaperSimSampler(H))).main() @ &m : res] +
    rejection_sampling_loss_term <=
    Pr[MLWE_Game(Bmlwe).main() @ &m : res] +
    Pr[BimodalSelfTargetMSIS_Game(Bbimodal).main() @ &m : res] =>
  Pr[MLWE_Game(Bmlwe).main() @ &m : res] <= mlwe_hardness_term =>
  Pr[ModuleSIS_Game(BimodalMSIS_As_ModuleSIS(Bbimodal)).main()
       @ &m : res] <= module_sis_hardness_term =>
  Pr[SIG.EUF_CMA(H, HAETAE, A).main() @ &m : res] <=
    haetae_euf_counted_rom_bound.
\end{lstlisting}
\noindent\textbf{Formal mathematical reading.}
Mathematical theorem. This theorem has the form \(\Pr[G:E] \le B\) is a game-probability statement.  It is one of the exact hops, bad-event bounds, or final advantage inequalities used in the reduction.

\noindent\textbf{Role in the proof.}
Use in this chapter. It contributes directly to an advantage bound, bad-event bound, or real-arithmetic composition step.

\noindent\textbf{Proof-script interpretation.}
Proof-script reading. The machine proof decomposes the theorem into rewrites, program-logic obligations, relational equivalences, and arithmetic side conditions.
\emph{Main tactics visible in the proof script:} \ec{apply}.
\emph{Named identifiers syntactically cited in the proof block:}
\begin{lstlisting}[language=EasyCrypt,basicstyle=\ttfamily\scriptsize]
euf_cma_security_from_ro_exact_hyperball_sampling_hop_and_counted_bimodal
lifted_loss
mlwe_bound
msis_bound
real_signing_ro_exact_hyperball_hop_from_paper_sample_lifting
ro_exact_nma_hop
sample_loss
\end{lstlisting}

\end{formaltheorem}

\begin{formaltheorem}[EasyCrypt line 1170]
\noindent\textbf{EasyCrypt name.}
\begin{lstlisting}[language=EasyCrypt,basicstyle=\ttfamily\scriptsize]
euf_cma_security_from_budgeted_paper_sample_lifting_and_counted_bimodal
\end{lstlisting}
\noindent\textbf{Checked EasyCrypt statement.}
\begin{lstlisting}[language=EasyCrypt,basicstyle=\ttfamily\scriptsize]
lemma euf_cma_security_from_budgeted_paper_sample_lifting_and_counted_bimodal
  &m :
  Pr[SIG.UF_NMA(H, HAETAE,
       ROMInternalTranscriptPaperSimAsNMA
         (A, ROSigningAttemptPaperSimSampler(H))).main() @ &m : res] <=
  Pr[SIG.UF_NMA(H, HAETAE,
       ROMInternalTranscriptBudgetedPaperSimAsNMA
         (A, ROSigningAttemptPaperSimSampler(H))).main() @ &m : res] =>
  Pr[SIG.UF_NMA(H, HAETAE,
       ROMInternalTranscriptBudgetedPaperSimAsNMA
         (A, ROExactHyperballPaperSimSampler(H))).main() @ &m : res] <=
  Pr[SIG.UF_NMA(H, HAETAE,
       ROMInternalTranscriptPaperSimAsNMA
         (A, ROExactHyperballPaperSimSampler(H))).main() @ &m : res] =>
  structural_to_exact_hyperball_paper_sample_loss_obligation haetae_mode =>
  Pr[SIG.UF_NMA(H, HAETAE,
       ROMInternalTranscriptPaperSimAsNMA
         (A, ROExactHyperballPaperSimSampler(H))).main() @ &m : res] +
    rom_signature_query_budget * rejection_sampling_loss_term <=
    Pr[MLWE_Game(Bmlwe).main() @ &m : res] +
    Pr[BimodalSelfTargetMSIS_Game(Bbimodal).main() @ &m : res] =>
  Pr[MLWE_Game(Bmlwe).main() @ &m : res] <= mlwe_hardness_term =>
  Pr[ModuleSIS_Game(BimodalMSIS_As_ModuleSIS(Bbimodal)).main()
       @ &m : res] <= module_sis_hardness_term =>
  Pr[SIG.EUF_CMA(H, HAETAE, A).main() @ &m : res] <=
    haetae_euf_counted_rom_bound.
\end{lstlisting}
\noindent\textbf{Formal mathematical reading.}
Mathematical theorem. This theorem has the form \(\Pr[G:E] \le B\) is a game-probability statement.  It is one of the exact hops, bad-event bounds, or final advantage inequalities used in the reduction.

\noindent\textbf{Role in the proof.}
Use in this chapter. It contributes directly to an advantage bound, bad-event bound, or real-arithmetic composition step.

\noindent\textbf{Proof-script interpretation.}
Proof-script reading. The machine proof decomposes the theorem into rewrites, program-logic obligations, relational equivalences, and arithmetic side conditions.
\emph{Main tactics visible in the proof script:} \ec{apply}.
\emph{Named identifiers syntactically cited in the proof block:}
\begin{lstlisting}[language=EasyCrypt,basicstyle=\ttfamily\scriptsize]
HAETAE
Pr
ROExactHyperballPaperSimSampler
ROMInternalTranscriptPaperSimAsNMA
SIG
UF_NMA
attempt_unbudgeted_to_budgeted
euf_cma_security_from_haetae_rejection_paper_sim_nma_and_counted_bimodal
exact_budgeted_to_unbudgeted
haetae_rejection_ro_exact_hyperball_sampling_hop_from_budgeted_lifting
ler_trans
main
mlwe_bound
msis_bound
rejection_sampling_loss_term
ro_exact_nma_hop
rom_signature_query_budget
sample_loss
\end{lstlisting}

\end{formaltheorem}

\begin{formaltheorem}[EasyCrypt line 1217]
\noindent\textbf{EasyCrypt name.}
\begin{lstlisting}[language=EasyCrypt,basicstyle=\ttfamily\scriptsize]
euf_cma_security_from_checked_ro_sampling_hop_and_counted_bimodal
\end{lstlisting}
\noindent\textbf{Checked EasyCrypt statement.}
\begin{lstlisting}[language=EasyCrypt,basicstyle=\ttfamily\scriptsize]
lemma euf_cma_security_from_checked_ro_sampling_hop_and_counted_bimodal &m :
  Pr[SIG.UF_NMA(H, HAETAE,
       ROMInternalTranscriptPaperSimAsNMA
         (A, ROExactHyperballPaperSimSampler(H))).main() @ &m : res] +
    rejection_sampling_loss_term <=
    Pr[MLWE_Game(Bmlwe).main() @ &m : res] +
    Pr[BimodalSelfTargetMSIS_Game(Bbimodal).main() @ &m : res] =>
  Pr[MLWE_Game(Bmlwe).main() @ &m : res] <= mlwe_hardness_term =>
  Pr[ModuleSIS_Game(BimodalMSIS_As_ModuleSIS(Bbimodal)).main()
       @ &m : res] <= module_sis_hardness_term =>
  Pr[SIG.EUF_CMA(H, HAETAE, A).main() @ &m : res] <=
    haetae_euf_counted_rom_bound.
\end{lstlisting}
\noindent\textbf{Formal mathematical reading.}
Mathematical theorem. This theorem has the form \(\Pr[G:E] \le B\) is a game-probability statement.  It is one of the exact hops, bad-event bounds, or final advantage inequalities used in the reduction.

\noindent\textbf{Role in the proof.}
Use in this chapter. It contributes directly to an advantage bound, bad-event bound, or real-arithmetic composition step.

\noindent\textbf{Proof-script interpretation.}
Proof-script reading. The machine proof decomposes the theorem into rewrites, program-logic obligations, relational equivalences, and arithmetic side conditions.
\emph{Main tactics visible in the proof script:} \ec{apply}.
\emph{Named identifiers syntactically cited in the proof block:}
\begin{lstlisting}[language=EasyCrypt,basicstyle=\ttfamily\scriptsize]
euf_cma_security_from_ro_exact_hyperball_sampling_hop_and_counted_bimodal
mlwe_bound
msis_bound
real_signing_ro_exact_hyperball_hop_checked_from_loss_budget
ro_exact_nma_hop
\end{lstlisting}

\end{formaltheorem}

\begin{formaltheorem}[EasyCrypt line 1239]
\noindent\textbf{EasyCrypt name.}
\begin{lstlisting}[language=EasyCrypt,basicstyle=\ttfamily\scriptsize]
euf_cma_security_from_real_signing_exact_hyperball_hop_and_counted_bimodal
\end{lstlisting}
\noindent\textbf{Checked EasyCrypt statement.}
\begin{lstlisting}[language=EasyCrypt,basicstyle=\ttfamily\scriptsize]
lemma euf_cma_security_from_real_signing_exact_hyperball_hop_and_counted_bimodal &m :
  Pr[SIG.UF_NMA(H, HAETAE,
       ROMInternalTranscriptPaperSimAsNMA
         (A, RealSigningPaperSimSampler(H))).main() @ &m : res] <=
  Pr[SIG.UF_NMA(H, HAETAE,
       ROMInternalTranscriptPaperSimAsNMA
         (A, ExactHyperballPaperSimSampler(H))).main() @ &m : res] +
    rejection_sampling_loss_term =>
  Pr[SIG.UF_NMA(H, HAETAE,
       ROMInternalTranscriptPaperSimAsNMA
         (A, ExactHyperballPaperSimSampler(H))).main() @ &m : res] +
    rejection_sampling_loss_term <=
    Pr[MLWE_Game(Bmlwe).main() @ &m : res] +
    Pr[BimodalSelfTargetMSIS_Game(Bbimodal).main() @ &m : res] =>
  Pr[MLWE_Game(Bmlwe).main() @ &m : res] <= mlwe_hardness_term =>
  Pr[ModuleSIS_Game(BimodalMSIS_As_ModuleSIS(Bbimodal)).main()
       @ &m : res] <= module_sis_hardness_term =>
  Pr[SIG.EUF_CMA(H, HAETAE, A).main() @ &m : res] <=
    haetae_euf_counted_rom_bound.
\end{lstlisting}
\noindent\textbf{Formal mathematical reading.}
Mathematical theorem. This theorem has the form \(\Pr[G:E] \le B\) is a game-probability statement.  It is one of the exact hops, bad-event bounds, or final advantage inequalities used in the reduction.

\noindent\textbf{Role in the proof.}
Use in this chapter. It contributes directly to an advantage bound, bad-event bound, or real-arithmetic composition step.

\noindent\textbf{Proof-script interpretation.}
Proof-script reading. The machine proof decomposes the theorem into rewrites, program-logic obligations, relational equivalences, and arithmetic side conditions.
\emph{Main tactics visible in the proof script:} \ec{apply}.
\emph{Named identifiers syntactically cited in the proof block:}
\begin{lstlisting}[language=EasyCrypt,basicstyle=\ttfamily\scriptsize]
euf_cma_security_from_exact_hyperball_sampling_hop_and_counted_bimodal
exact_nma_hop
haetae_rejection_exact_hyperball_sampling_hop_from_real_signing_exact_hyperball_hop
mlwe_bound
msis_bound
real_exact_hop
\end{lstlisting}

\end{formaltheorem}

\subsubsection*{Explicit Program-Logic Judgments Used Inside Proof Scripts}
\begin{formaljudgment}[EasyCrypt line 44]
\noindent\textbf{Checked EasyCrypt judgment.}
\begin{lstlisting}[language=EasyCrypt,basicstyle=\ttfamily\scriptsize]
hoare.
\end{lstlisting}
\noindent\textbf{Formal mathematical reading.} This is a Hoare judgment.  It states partial correctness of the displayed procedure from the displayed precondition to the displayed postcondition.
\end{formaljudgment}

\medskip
\noindent\emph{End of generated formal inventory for \file{HAETAE_Security.ec}.}
